% Official Preprints.org class, 23 June 2026; see TEMPLATE_SOURCE.md.
\documentclass[preprints,article,accept,oneauthor]{Definitions/mdpi}
\setlength{\headheight}{18pt}

\firstpage{1}
\makeatletter
\setcounter{page}{\@firstpage}
\makeatother
\pubvolume{1}
\issuenum{1}
\articlenumber{0}
\pubyear{2026}
\copyrightyear{2026}
\datereceived{}
\daterevised{}
\dateaccepted{}
\datepublished{}

% amsmath, amssymb, amsthm, natbib, hyperref and cleveref are class-supplied.
\usepackage{mathtools}
\usepackage{placeins}
\allowdisplaybreaks
\numberwithin{equation}{section}

% The class reserves counters named theorem, lemma, and proposition.
\newcounter{result}[section]
\renewcommand{\theresult}{\thesection.\arabic{result}}
\theoremstyle{mdpi}
\newtheorem{theorem}[result]{Theorem}
\newtheorem{lemma}[result]{Lemma}
\newtheorem{proposition}[result]{Proposition}
\newtheorem{corollary}[result]{Corollary}
\theoremstyle{definition}
\newtheorem{remark}[result]{Remark}

\DeclareMathOperator{\diag}{diag}
\DeclareMathOperator{\tridiag}{tridiag}
\newcommand{\C}{\mathbb C}
\newcommand{\R}{\mathbb R}
\newcommand{\ii}{\mathrm i}
\newcommand{\proofstage}[1]{%
  \par\medskip\noindent\textit{#1}\par\nobreak\smallskip}

\Title{Connectedness thresholds for pseudospectra of tridiagonal Toeplitz matrices}
\newcommand{\orcidauthorA}{0009-0009-5346-753X}
\Author{Shanmu Jin\orcidA{}}
\AuthorNames{Shanmu Jin}
\address[1]{Department of Neurosurgery, Peking Union Medical College Hospital,
No.~1 Shuaifuyuan, Dongcheng District, Beijing 100730, China;\\
Email: \href{mailto:jinshanmu1129@pku.edu.cn}{jinshanmu1129@pku.edu.cn}}
\corres{Correspondence: jinshanmu1129@pku.edu.cn}

\abstract{\texorpdfstring{We study the open Euclidean-norm pseudospectra of
\(A_n(a)=\operatorname{tridiag}(a,0,1)\), where \(0<a<1\).
For each order \(n\), connectedness is equivalent to
\(\varepsilon>\gamma_n(a)\), where \(\gamma_n(a)\) is the maximum of
\(s_{\min}(xI_n-A_n(a))\) over the real spectral interval.
Every pseudospectral component is contractible.  The thresholds decrease
strictly with \(n\), so connectedness, once attained, persists at every
larger order.  The proof uses vertical monotonicity of the least singular
value, a symmetric matrix pencil, and separate comparisons for noncentral
and central spectral gaps.  Explicit upper and lower bounds show that
\(\gamma_n(a)\asymp_a{(n+1)a^{n/2}}\).  Consequently, the first connected
order equals
\(2[\log(1/\varepsilon)+\log\log(1/\varepsilon)]/\log(1/a)+O_a(1)\)
as \(\varepsilon\) tends to zero.  We also give the connectedness thresholds
at \(a=0\) and \(a=1\).}{Exact connectedness thresholds, strict decrease with matrix order, and quantitative bounds for the Euclidean-norm pseudospectra of tridiagonal Toeplitz matrices.}}

\keyword{pseudospectrum; tridiagonal Toeplitz matrix; least singular value;
connectedness}
\MSC{15A18; 15A60; 15B05; 65F15}


\begin{document}
\crefname{theorem}{Theorem}{Theorems}
\Crefname{theorem}{Theorem}{Theorems}
\crefname{lemma}{Lemma}{Lemmas}
\Crefname{lemma}{Lemma}{Lemmas}
\crefname{proposition}{Proposition}{Propositions}
\Crefname{proposition}{Proposition}{Propositions}
\crefname{section}{Section}{Sections}
\Crefname{section}{Section}{Sections}

\section{Introduction}
\label{sec:introduction}

For a nonnormal matrix, small perturbations can move eigenvalues far from
its spectrum.  The pseudospectrum records this sensitivity, and its
connected components describe which spectral clusters remain separated
at a prescribed perturbation level.  This paper determines when the open
Euclidean-norm pseudospectrum of
\begin{equation}\label{eq:canonical-matrix}
 A_n(a)=\tridiag(a,0,1)
 =\begin{bmatrix}
 0&1\\
 a&0&\ddots\\
 &\ddots&\ddots&1\\
 &&a&0
 \end{bmatrix},\qquad n\geq2,\quad 0<a<1,
\end{equation}
is connected, and proves that connectedness persists as the order increases.
The arguments of \(\tridiag\) specify the subdiagonal, diagonal, and
superdiagonal, respectively.

Throughout, \(\|\cdot\|_2\) denotes the Euclidean vector norm or its induced
matrix norm, \(s_{\min}\) the least singular value, and \(I_n\) the identity
matrix of order \(n\).  We use the open pseudospectrum
\begin{equation}\label{eq:intro-pseudospectrum}
 \sigma_\varepsilon(A_n)
 =\{z\in\C:s_{\min}(zI_n-A_n)<\varepsilon\},\qquad\varepsilon>0.
\end{equation}
It is also the union of \(\sigma(A_n+E)\) over complex matrices
\(E\) with \(\|E\|_2<\varepsilon\), or, equivalently, the spectrum together
with the set where \(\|(zI_n-A_n)^{-1}\|_2>\varepsilon^{-1}\)
\cite{Trefethen1997}.  Thus the perturbations in this paper are unrestricted
complex matrices, measured in the induced Euclidean norm.

Toeplitz pseudospectra have long provided examples of the difference between
spectral location and perturbation sensitivity.  Resolvent estimates and
pseudospectral behavior of finite non-Hermitian Toeplitz matrices are studied
in \cite{ReichelTrefethen1992}.  Limiting results for pseudospectra and singular
values of large truncated Wiener--Hopf operators are developed in \cite{Bottcher1994}.  Explicit eigensystems and sensitivity
formulas for tridiagonal Toeplitz matrices are given in
\cite{NoschesePasquiniReichel2013}.  Algorithms for the extreme points of
Toeplitz-structured pseudospectra in the Frobenius norm are developed in
\cite{ButtaGuglielmiNoschese2012}.  Finite-section methods also provide
convergent spectral and pseudospectral inclusion sets for banded and
band-dominated operators \cite{ChandlerWildeEtAl2024}.

The geometry of pseudospectra is closely connected with matrix nearness
problems.  Their component structure and variational properties are studied
in \cite{BurkeLewisOverton2003}.  For a matrix with distinct eigenvalues,
the first coalescence level of pseudospectral components equals the distance to matrices with a multiple eigenvalue
\cite{BurkeLewisOverton2007}.  First-coalescence points are lowest generalized saddle points of the
least-singular-value function \cite{AlamEtAl2011}.  Here the question is when all components have joined,
and how that level changes between consecutive matrix orders.

Put \(r=\sqrt a\), and let \(V_n=\tridiag(0,0,1)\).  The diagonal matrix
\(\mathcal G_n(r)=\diag(1,r,\ldots,r^{n-1})\) gives
\begin{equation}\label{eq:intro-similarity}
 A_n=\mathcal G_n(r)\,r(V_n+V_n^T)\,\mathcal G_n(r)^{-1}.
\end{equation}
The spectrum is therefore
\begin{equation}\label{eq:eigenvalues}
 \lambda_{n,k}=2r\cos\frac{k\pi}{n+1},\qquad 1\leq k\leq n.
\end{equation}
Write
\begin{equation}\label{eq:barrier-definition}
 \begin{aligned}
 I_n^{\mathrm{sp}}&=[\lambda_{n,n},\lambda_{n,1}],&
 g_n(x)&=s_{\min}(xI_n-A_n),\\
 \gamma_n(a)&=\max_{x\in I_n^{\mathrm{sp}}}g_n(x),&
 N_c(a,\varepsilon)&=\min\{n\geq2:
 \sigma_\varepsilon(A_n(a))\text{ is connected}\}.
 \end{aligned}
\end{equation}
The maximum exists by continuity.  The following theorem also ensures that
\(N_c\) is finite.  When \(a\) is fixed, we suppress it in \(A_n\),
\(g_n\), and \(\gamma_n\).  We set \(A_1=V_1=[0]\) when a smaller block
occurs in a proof.

\begin{theorem}\label[theorem]{thm:canonical-main}
Let \(0<a<1\).  For every \(n\geq2\) and \(\varepsilon>0\), each connected
component of \(\sigma_\varepsilon(A_n(a))\) is contractible, and
\begin{equation}\label{eq:connected-iff-gamma}
 \sigma_\varepsilon(A_n(a))\text{ is connected}
 \quad\Longleftrightarrow\quad \varepsilon>\gamma_n(a).
\end{equation}
Moreover,
\begin{equation}\label{eq:main-strict}
 \gamma_{n+1}(a)<\gamma_n(a),\qquad
 \gamma_2(a)=a,\qquad \lim_{n\to\infty}\gamma_n(a)=0.
\end{equation}
Consequently,
\[
 \{n\geq2:\sigma_\varepsilon(A_n(a))\text{ is connected}\}
 =\{n\geq N_c(a,\varepsilon)\},
\]
and \(N_c(a,\varepsilon)=2\) if and only if \(\varepsilon>a\).
For fixed \(a\),
\begin{equation}\label{eq:N-asymptotic-main}
 N_c(a,\varepsilon)=\frac{2}{\log(1/a)}
 \left(\log(1/\varepsilon)+\log\!\log(1/\varepsilon)\right)+O_a(1)
 \qquad(\varepsilon\downarrow0).
\end{equation}
\end{theorem}

The criterion \eqref{eq:connected-iff-gamma} follows from contraction toward
the real axis.  The main comparison is \(\gamma_{n+1}<\gamma_n\): both the
spectral intervals and their gaps change with \(n\), and the central gaps
change with parity.  We compare noncentral gaps through a parametrization
of a symmetric matrix pencil, then treat the central gaps separately.
Finally, explicit bounds
\(\gamma_n(a)\asymp_a(n+1)a^{n/2}\) give
\eqref{eq:N-asymptotic-main}.

Section~\ref{sec:vertical} proves contraction and the connectedness criterion.
Section~\ref{sec:mesh} establishes the strict comparison, including the
continuant identities used in its proof.  Section~\ref{sec:topology-bounds}
gives quantitative bounds and proves the remaining assertions of the theorem.
The boundary cases \(a=0,1\) and numerical illustrations follow in
Sections~\ref{sec:boundaries} and~\ref{sec:numerics}.


\section{Contraction toward the real axis}
\label{sec:vertical}

The least singular value is nondecreasing along every vertical half-line
issuing from the real axis.  We prove this through a determinant identity
and then use it to describe the pseudospectral components.

\subsection{A determinant identity}

\begin{proposition}[Positivity of the determinant derivative]
\label[proposition]{prop:vertical-determinant-certificate}
Fix \(0<a<1\), \(x\in\mathbb R\), \(Y\geq0\), and \(t\in\mathbb R\).  For \(k\geq1\), set
\[
 \begin{aligned}
 \mathcal P_k(Y,t)
 &=((x+\ii\sqrt Y)I_k-A_k)^*((x+\ii\sqrt Y)I_k-A_k)-t^2I_k,\\
 F_k(Y,t)&=\det\mathcal P_k(Y,t),
 \end{aligned}
\]
and set \(F_0=1\).  For every \(n\geq1\),
\begin{equation}\label{eq:det-derivative-positive}
 \mathcal P_n(Y,t)>0
 \quad\Longrightarrow\quad
 \partial_YF_n(Y,t)>0.
\end{equation}
\end{proposition}

\begin{proof}
At a fixed value of \(Y\), write
\[
 \delta_{\rm pent}=x^2+Y+1+a^2-t^2,
 \qquad b=-(1+a)x+\ii(1-a)\sqrt Y.
\]
For \(n\geq2\), the matrix \(\mathcal P_n(Y,t)\) has diagonal
\[
 \delta_{\rm pent}-1,\delta_{\rm pent},\ldots,\delta_{\rm pent},
 \delta_{\rm pent}-a^2,
\]
whereas for \(n=1\) its sole entry is
\(\delta_{\rm pent}-1-a^2\).  In all dimensions
its first upper diagonal is \(b\), its first lower diagonal is \(\bar b\),
and its second off-diagonals are \(a\).  Let \(T_k\) be the determinant of
the pure pentadiagonal Toeplitz matrix with diagonal
\(\delta_{\rm pent}\) and these two
off-diagonals.  Set \(T_{-1}=0,T_0=1\), and put
\begin{equation}\label{eq:D-R-def}
 D_k=T_k-T_{k-1},\qquad R_k=T_k-a^2T_{k-1}.
\end{equation}
Since a determinant is affine in each diagonal entry, replacing the
first diagonal entry \(\delta_{\rm pent}\) by
\(\delta_{\rm pent}-1\) subtracts the cofactor \(T_{k-1}\), whereas
replacing the last one by \(\delta_{\rm pent}-a^2\) subtracts
\(a^2T_{k-1}\).  Thus, for \(k<n\), \(D_k\) and \(R_k\) are,
respectively, the leading and trailing principal minors of
\(\mathcal P_n(Y,t)\).

The following continuant calculation isolates all dependence on \(Y\).
With \(T_k=0\) for \(k<0\), a Laplace expansion over the last two rows,
grouped according to the largest displaced index in the contributing
permutation, gives for every \(k\geq1\)
\begin{equation}\label{eq:T-recurrence}
\begin{aligned}
T_k={}&\delta_{\rm pent}T_{k-1}-(|b|^2-a^2)T_{k-2}
-\bigl(2a^2\delta_{\rm pent}-a(b^2+\bar b^2)\bigr)T_{k-3}\\
&-a^2(|b|^2-a^2)T_{k-4}+a^4\delta_{\rm pent}T_{k-5}-a^6T_{k-6}
-a^2\boldsymbol 1_{\{k=2\}}.
\end{aligned}
\end{equation}
For example, \(T_1=\delta_{\rm pent}\) and
\(T_2=\delta_{\rm pent}^2-|b|^2\); these also verify the
only inhomogeneous coefficient in \eqref{eq:T-recurrence}.  Multiplying
the recurrence by \(z^k\) and summing gives
\begin{equation}\label{eq:T-generating}
 \sum_{k\geq0}T_kz^k=\frac{1-a^2z^2}{Q(z,Y)},
\end{equation}
where
\[
\begin{aligned}
 Q(z,Y)={}&1-\delta_{\rm pent}z+(|b|^2-a^2)z^2
 +\bigl(2a^2\delta_{\rm pent}-a(b^2+\bar b^2)\bigr)z^3\\
 &+a^2(|b|^2-a^2)z^4-a^4\delta_{\rm pent}z^5+a^6z^6.
\end{aligned}
\]
Substituting \eqref{eq:T-generating} into \eqref{eq:D-R-def}, with
\(\mathcal D(z,Y)=\sum_{k\geq0}D_kz^k\), gives the first identity in
\begin{equation}\label{eq:D-F-generating}
 \mathcal D(z,Y)=\frac{(1-z)(1-a^2z^2)}{Q(z,Y)},\qquad
 \sum_{k\geq0}F_k(Y,t)z^k=(1-a^2z)\mathcal D(z,Y).
\end{equation}
For the second identity in \eqref{eq:D-F-generating}, first change the initial diagonal, producing
\(D_k\), and then change the final diagonal; its cofactor is
\(D_{k-1}\).  Thus \(F_k=D_k-a^2D_{k-1}\), which proves the asserted
generating function.  Direct substitution of \(\delta_{\rm pent}\) and \(b\) gives
\[
 \partial_YQ=-z(1-z)(1+az)^2(1-a^2z).
\]
Differentiating \eqref{eq:D-F-generating} therefore yields
\begin{equation}\label{eq:F-prime-generating}
 \sum_{k\geq0}\partial_YF_k(Y,t)z^k
 =\varphi_a(z)\mathcal D(z,Y)^2,
 \qquad
 \varphi_a(z)=\frac{z(1+az)(1-a^2z)^2}{1-az}.
\end{equation}

Suppose that \(\mathcal P_n(Y,t)>0\).  Then
\begin{equation}\label{eq:D-R-positive}
 D_k>0,\qquad R_k>0\qquad(0\leq k<n).
\end{equation}
Under the positivity in \eqref{eq:D-R-positive}, put
\(u_k=\sum_{j=0}^kD_jD_{k-j}\), with \(u_k=0\) for \(k<0\).
To extract coefficients from \eqref{eq:F-prime-generating}, note that the
coefficients of \(\varphi_a\), beginning with the coefficient of
\(z\), are
\[
 1,\quad 2a(1-a),\quad a^2(a^2-4a+2),\quad
 2a^{k-1}(1-a)^2\quad(k\geq4).
\]
Thus
\begin{equation}\label{eq:F-prime-convolution}
\begin{aligned}
 \partial_YF_n={}&u_{n-1}+2a(1-a)u_{n-2}
 +a^2(a^2-4a+2)u_{n-3}\\
 &+2(1-a)^2\sum_{k=4}^na^{k-1}u_{n-k}.
\end{aligned}
\end{equation}

The third coefficient is the sole term that may be negative; the following
cofactor inequality controls its contribution.  Let \(j,k\geq1\),
\(j+k=n-1\).  Deleting index \(k+1\) from \(\mathcal P_n\) leaves blocks of
sizes \(k\) and \(j\), with the second off-diagonal \(a\) as their sole
coupling.  In block form, the deleted principal submatrix is
\[
 \begin{bmatrix}
  L_k&a \mathbf e_k\mathbf e_1^T\\
  a \mathbf e_1\mathbf e_k^T&R_j^{\rm mat}
 \end{bmatrix},
 \qquad \det L_k=D_k,\quad \det R_j^{\rm mat}=R_j.
\]
Expanding once in each of the two coupling entries shows that its
determinant is
\begin{equation}\label{eq:deleted-cofactor}
 D_kR_j-a^2D_{k-1}R_{j-1}>0.
\end{equation}
Deleting index \(j+1\) instead gives the same inequality with \(j,k\)
interchanged.  Define
\[
 \vartheta_k=\frac{D_k}{D_{k-1}},\qquad
 \varpi_k=\frac{R_k}{R_{k-1}}.
\]
Since \(T_k=\sum_{\ell=0}^kD_\ell\), \eqref{eq:D-R-def} gives
\begin{equation}\label{eq:R-D-relation}
 R_k=D_k+(1-a^2)\sum_{\ell=0}^{k-1}D_\ell,\qquad
 R_k-R_{k-1}=D_k-a^2D_{k-1}.
\end{equation}
The second identity in \eqref{eq:R-D-relation} shows that
\(\vartheta_k\leq a^2\) implies \(\varpi_k\leq1\).  On dividing
\eqref{eq:deleted-cofactor} and its companion by positive factors, we
obtain
\begin{equation}\label{eq:ratio-cross}
 \vartheta_k\varpi_j>a^2,\qquad \vartheta_j\varpi_k>a^2.
\end{equation}
If both \(\vartheta_j\) and \(\vartheta_k\) were at most \(a^2\), then
both \(\varpi_j\) and \(\varpi_k\) would be at most \(1\),
contradicting \eqref{eq:ratio-cross}.  Hence, after interchanging
\(j,k\) if necessary, \(\vartheta_j>a^2\).  Moreover,
\[
 \vartheta_jR_{j-1}-a^2R_j
 =(\vartheta_j-a^2)(R_{j-1}-a^2D_{j-1}),
\]
and the second factor equals
\((1-a^2)\sum_{\ell=0}^{j-1}D_\ell>0\).  Therefore
\[
 \vartheta_j>a^2\quad\Longrightarrow\quad
 \varpi_j<\vartheta_j/a^2.
\]
The first inequality in \eqref{eq:ratio-cross} now gives
\(a^2<\vartheta_k\varpi_j<\vartheta_j\vartheta_k/a^2\), and therefore
\begin{equation}\label{eq:D-ratio-product}
 \vartheta_j\vartheta_k>a^4.
\end{equation}
Consequently, for \(n\geq3\),
applying \eqref{eq:D-ratio-product} to \(j=p+1\), \(k=n-2-p\) gives
\[
 D_{p+1}D_{n-2-p}>a^4D_pD_{n-3-p},
 \qquad0\leq p\leq n-3.
\]
Summing and then adding the two positive endpoint terms gives
\begin{equation}\label{eq:u-absorption}
 u_{n-1}=2D_{n-1}+\sum_{j=1}^{n-2}D_jD_{n-1-j}>a^4u_{n-3}.
\end{equation}

If the third coefficient in \eqref{eq:F-prime-convolution} is negative,
write it as \(-\mathcal B_{\rm neg}\), where
\(\mathcal B_{\rm neg}=a^2(4a-2-a^2)\).  Since
\[
 a^4-\mathcal B_{\rm neg}=2a^2(1-a)^2>0,
\]
equations \eqref{eq:F-prime-convolution} and \eqref{eq:u-absorption}
prove \eqref{eq:det-derivative-positive}.  If that coefficient is
nonnegative, the implication follows directly.  For \(n=1\), one has
\(\partial_YF_1=1\); for \(n=2\), one has
\(\partial_YF_2=2D_1+2a(1-a)>0\).  Thus
\eqref{eq:det-derivative-positive} holds in every dimension, which proves
the proposition.
\end{proof}


\subsection{Vertical monotonicity}

\begin{lemma}[Vertical monotonicity]\label[lemma]{lem:vertical}
Let \(n\geq2\) and \(0<a<1\).  For every real \(x\),
\[
 y\longmapsto s_{\min}((x+\ii y)I_n-A_n),\qquad y\geq0,
\]
is nondecreasing.
\end{lemma}

\begin{proof}
Use the notation of \Cref{prop:vertical-determinant-certificate} and put
\(Y=y^2\).

Fix \(Y_0\geq0\).  When the least singular value at \(Y_0\) equals
zero, monotonicity follows from nonnegativity of singular values.  Otherwise
take \(0\leq t<s_{\min}((x+\ii\sqrt{Y_0})I_n-A_n)\).  On the maximal interval
to the right of \(Y_0\) on which \(\mathcal P_n(Y,t)>0\), implication
\eqref{eq:det-derivative-positive} makes its determinant strictly increasing.
At any finite right endpoint, continuity of Hermitian eigenvalues would
produce a positive-semidefinite singular matrix with determinant zero.
Strict increase from the positive determinant at \(Y_0\) precludes such an
endpoint.
Thus \(\mathcal P_n(Y,t)>0\) for every \(Y\geq Y_0\).  Letting \(t\) increase to
the least singular value at \(Y_0\) proves the lemma.
\end{proof}
\subsection{Components and the real-axis criterion}

We first recall the component property of finite-dimensional
pseudospectra \cite{BurkeLewisOverton2003} and include its short proof.

\begin{lemma}\label[lemma]{lem:component}
Every connected component of \(\sigma_\varepsilon(A_n)\) contains an
 eigenvalue of \(A_n\).
\end{lemma}

\begin{proof}
The estimate \(s_{\min}(zI_n-A_n)\geq |z|-\|A_n\|_2\) makes the
pseudospectrum bounded.  Suppose that a component \(\mathcal C\) contains
no eigenvalue.  Its closure contains no eigenvalue either: a sufficiently
small pseudospectral disk about such an eigenvalue would meet
\(\mathcal C\) and hence belong to the same component.  Thus the resolvent
is analytic on a neighborhood of \(\overline{\mathcal C}\).
At every boundary point its norm is at most \(\varepsilon^{-1}\).
Indeed, a boundary point in the open pseudospectrum would have a small
connected neighborhood in that set and would belong to \(\mathcal C\).
The resolvent norm is subharmonic, so the maximum principle bounds it by
\(\varepsilon^{-1}\) throughout \(\mathcal C\), a contradiction.
\end{proof}

\begin{proposition}\label[proposition]{prop:topology}
Each component of \(\sigma_\varepsilon(A_n)\) is contractible, and
\eqref{eq:connected-iff-gamma} holds.
\end{proposition}

\begin{proof}
Since \(A_n\) is real, conjugation preserves the least singular value.
Lemma~\ref{lem:vertical} therefore shows that
\[
 H_t(x+\ii y)=x+\ii(1-t)y,\qquad0\leq t\leq1,
\]
maps \(\Omega=\sigma_\varepsilon(A_n)\) into itself.  For any component
\(\mathcal C\), the path \(t\mapsto H_t(z)\) remains in \(\mathcal C\).
The restriction of \(H_t\) is a strong deformation retraction onto
\(\mathcal C\cap\R\).  This intersection is nonempty and connected,
being the image of \(\mathcal C\) under \(H_1\); hence it is an interval.
It follows that \(\mathcal C\) is contractible.

If \(\Omega\) is connected, its projection \(\Omega\cap\R\) is an interval
containing all eigenvalues, and therefore contains \(I_n^{\mathrm{sp}}\).
Conversely, if \(I_n^{\mathrm{sp}}\subset\Omega\), that interval belongs
to one component and contains every eigenvalue.  Lemma~\ref{lem:component}
excludes any other component.  Thus connectedness is equivalent to
\(g_n(x)<\varepsilon\) on the compact interval \(I_n^{\mathrm{sp}}\),
or \(\gamma_n<\varepsilon\).  At equality a maximizing point is excluded,
so the open pseudospectrum is disconnected.
\end{proof}


\section{Strict matrix-order comparison of real-gap barriers}
\label{sec:mesh}

The global barrier may occur in any inter-eigenvalue gap, while the central
gap changes with parity.  We therefore compare every gap.  Persymmetry
selects the signed eigenvalue branch whose magnitude is the least singular
value; folded Chebyshev coordinates and chord geometry compare the
noncentral gaps, and a separate argument treats the parity-changing central
gaps.  Write the eigenvalues in decreasing order as
\[
 \lambda_{n,k}=2r\cos\frac{k\pi}{n+1},\qquad 1\leq k\leq n,
\]
and call
\[
 G_{n,k}=[\lambda_{n,k+1},\lambda_{n,k}],\qquad
 \gamma_{n,k}=\max_{x\in G_{n,k}}g_n(x),\qquad 1\leq k<n,
\]
the closed spectral gaps and their heights.  Continuity guarantees that each
maximum is attained.  The spectral endpoints have height zero, whereas
\(g_n(x)>0\) in every open gap.  Thus every gap height is positive, every
maximizer lies in the open gap, and
\begin{equation}\label{eq:gamma-def}
 \gamma_n=\max_{1\leq k<n}\gamma_{n,k}.
\end{equation}
\begin{lemma}[Strict barrier comparison]\label[lemma]{lem:mesh}
For every \(n\geq2\),
\[
 \gamma_{n+1}<\gamma_n.
\]
\end{lemma}

\subsection{Signed-pencil representation and branch selection}
\label{subsec:signed-pencil}

Let \(J_n\) be the reversal permutation matrix, so that
\(J_n\mathbf e_j=\mathbf e_{n+1-j}\).
Persymmetry supplies the signed spectral representation needed for the
matrix-order comparison:
\(A_nJ_n=J_nA_n^T\), so
\(B_n(x)=(xI_n-A_n)J_n\) is real symmetric.  Moreover,
\[
 (xI_n-A_n-sJ_n)J_n=B_n(x)-sI_n.
\]
Since \(J_n\) is orthogonal, the absolute values of the eigenvalues of \(B_n(x)\)
are exactly the singular values of \(xI_n-A_n\), including multiplicity.
Consequently,
\begin{equation}\label{eq:signed-pencil}
 s_{\min}(xI_n-A_n)
 =\min\{|s|:\det(xI_n-A_n-sJ_n)=0\}.
\end{equation}
Set \(q_n(x,s)=\det(xI_n-A_n-sJ_n)\); the minimum in
\eqref{eq:signed-pencil} selects its signed branch nearest zero.
\begin{lemma}[The middle signed branch]\label[lemma]{lem:middle-branch}
Write the eigenvalues of \(B_n(x)\) in decreasing order as
\(\mu_1(x)\geq\cdots\geq\mu_n(x)\), and abbreviate
\(\lambda_k=\lambda_{n,k}\).
If \(\lambda_{j+1}<x<\lambda_j\), then
\[
 s_{\min}(xI_n-A_n)=|s_*(x)|,
 \qquad \operatorname{sgn}(s_*(x))=(-1)^j,
\]
where
\[
 s_*(x)=
 \begin{cases}
  \mu_m(x),&n=2m,\\
  \mu_{m+1}(x),&n=2m+1
 \end{cases}.
\]
\end{lemma}

\begin{proof}
Let \(A_n^{\rm sym}=r(V_n+V_n^T)\).  The identities
\[
 \begin{aligned}
 A_n&=\mathcal G_n(r)A_n^{\rm sym}\mathcal G_n(r)^{-1},\\
 \mathcal G_n(r)^{-1}J_n
 &=r^{-(n-1)}J_n\mathcal G_n(r)
 \end{aligned}
\]
give
\begin{equation}\label{eq:B-congruence}
 B_n(x)=r^{-(n-1)}\mathcal G_n(r)(xI_n-A_n^{\rm sym})
 J_n\mathcal G_n(r).
\end{equation}
The sine eigenvector of \(A_n^{\rm sym}\) belonging to \(\lambda_k\) is an
eigenvector of \(J_n\) with eigenvalue \((-1)^{k+1}\).  Hence
\eqref{eq:B-congruence} and Sylvester's law of inertia show that
\(B_n(x)\) has the same inertia as the diagonal matrix with entries
\[
 (-1)^{k+1}(x-\lambda_k),\qquad 1\leq k\leq n.
\]
For a real symmetric matrix \(H\), let \(\nu_+(H)\) denote the number of its
positive eigenvalues, counted with multiplicity.  Counting signs gives
\begin{equation}\label{eq:middle-inertia}
 \nu_+(B_{2m}(x))=
 \begin{cases}m,&j\text{ even},\\m-1,&j\text{ odd},\end{cases}
 \qquad
 \nu_+(B_{2m+1}(x))=
 \begin{cases}m+1,&j\text{ even},\\m,&j\text{ odd}.
 \end{cases}
\end{equation}

First let \(n=2m\).  Permute the coordinates into the two groups
\[
 (1,\ldots,m),\qquad(2m,2m-1,\ldots,m+1).
\]
Here \(\mathbf e_j\) denotes the \(j\)th Euclidean coordinate vector in
the ambient space.  Put \(P_m(x)=xI_m-A_m\).  In these coordinates
\begin{equation}\label{eq:even-dilation}
\begin{aligned}
 B_{2m}(x)&=\widehat P_m(x)-\mathcal R_{\rm bd},\\
 \widehat P_m(x)&=\begin{bmatrix}0&P_m(x)\\P_m(x)^T&0\end{bmatrix},\\
 \mathcal R_{\rm bd}
 &=\begin{bmatrix}\mathbf e_m\mathbf e_m^T&0\\
 0&a \mathbf e_m\mathbf e_m^T\end{bmatrix},
\end{aligned}
\end{equation}
In \eqref{eq:even-dilation},
\(\mathcal R_{\rm bd}\geq0\) and
\(\operatorname{rank}\mathcal R_{\rm bd}=2\).  For
\(\mathsf S=\diag(I_m,-I_m)\),
\[
 \widetilde B=-\mathsf S B_{2m}(x)\mathsf S
 =\widehat P_m(x)+\mathcal R_{\rm bd}
 =B_{2m}(x)+2\mathcal R_{\rm bd}.
\]
If \(\widetilde\mu_1\geq\cdots\geq\widetilde\mu_{2m}\) are the
eigenvalues of
\(\widetilde B\), then
\begin{equation}\label{eq:alpha-reflection}
 \widetilde\mu_k=-\mu_{2m+1-k}.
\end{equation}
Combining the reflection \eqref{eq:alpha-reflection} with
positive-semidefinite rank-two interlacing gives
\begin{equation}\label{eq:rank-two-interlace}
 \widetilde\mu_k\geq\mu_k\geq\widetilde\mu_{k+2},
 \qquad1\leq k\leq2m-2.
\end{equation}
If \(j\) is even, the only candidates for minimum modulus
have moduli \(\mu_m\) and
\(-\mu_{m+1}=\widetilde\mu_m\), and
\(\mu_m\leq\widetilde\mu_m\).  If \(j\) is odd and \(m\geq2\), the two
candidates have moduli \(\mu_{m-1}\) and
\(-\mu_m=\widetilde\mu_{m+1}\); taking \(k=m-1\) in
\eqref{eq:rank-two-interlace} gives
\(\mu_{m-1}\geq\widetilde\mu_{m+1}\).  For \(m=1\), both eigenvalues in
the sole gap are negative, so the larger one \(\mu_1\) has minimum
modulus.  Thus \(\mu_m\) always attains the minimum modulus, and its
sign follows from \eqref{eq:middle-inertia}.  Equality in an interlacing
comparison corresponds to a multiple least singular value.

Now let \(n=2m+1\).  Group the first and last \(m\) coordinates as
above and leave the middle coordinate last.  Then
\begin{equation}\label{eq:odd-dilation}
 B_{2m+1}(x)=
 \begin{bmatrix}\widehat P_m(x)&\boldsymbol w\\\boldsymbol w^T&x\end{bmatrix},
 \qquad
 \boldsymbol w=\begin{bmatrix}-\mathbf e_m\\-a\mathbf e_m\end{bmatrix}.
\end{equation}
If \(s_1(P_m)\geq\cdots\geq s_m(P_m)\) are the singular values of
\(P_m(x)\), then the eigenvalues of the principal submatrix
\(\widehat P_m(x)\) are
\[
 s_1(P_m)\geq\cdots\geq s_m(P_m)\geq0\geq
 -s_m(P_m)\geq\cdots\geq-s_1(P_m).
\]
Cauchy interlacing gives
\[
 \mu_m\geq s_m(P_m)\geq\mu_{m+1}
 \geq-s_m(P_m)\geq\mu_{m+2}.
\]
If \(j\) is even, \eqref{eq:middle-inertia} and interlacing give
\[
 0<\mu_{m+1}\leq s_m(P_m)\leq-\mu_{m+2};
\]
if \(j\) is odd, they give
\[
 0<-\mu_{m+1}\leq s_m(P_m)\leq\mu_m.
\]
Every other eigenvalue is farther from zero.  Hence
\(\mu_{m+1}\) attains the minimum modulus, with sign prescribed by
\eqref{eq:middle-inertia}; equality corresponds to multiplicity.  At a
spectral endpoint the zero is simple because \(A_n\) is similar to the
Jacobi matrix \(A_n^{\rm sym}\), whose eigenvalues are simple.  Thus the selected
ordered branch extends continuously to zero at both ends of the closed
gap.  This proves the assertion.
\end{proof}

\subsection{Folded continuants}
\label{subsec:folded-continuants}

We use folded coordinates for the noncentral gaps and a separate
interlacing argument for the parity-changing central gaps.
Let \(U_j\) be the Chebyshev polynomial of the second kind, with
\(U_{-1}=0\).  Put
\[
 C_0=x^2+(1-a)^2-s^2,\qquad
 \zeta_s=\frac{1+a^2+(1+a)s}{2a},
\]
and define the folded discriminant and recurrence coefficient by
\[
 \Delta_{\rm fold}=((1+a-x)^2-s^2)((1+a+x)^2-s^2),
\]
\[
 \mathfrak b=(1+a^2)x^2+2as^2+2a^2-2a-2a^3.
\]
The global recurrence polynomial is
\begin{equation}\label{eq:fold-characteristic}
\begin{gathered}
 \Pi(z):=1-C_0z+\mathfrak b z^2-a^2C_0z^3+a^4z^4,\\
 u,v\in\mathbb R,\qquad
 2a(u+v)=C_0,\qquad a^2(4uv+2)=\mathfrak b,\\
 \Longrightarrow
 \Pi(z)=(1-2auz+a^2z^2)(1-2avz+a^2z^2).
\end{gathered}
\end{equation}
On a real sheet with \(\Delta_{\rm fold}\geq0\), we additionally put
\[
 \xi_\pm=\frac{C_0\pm\sqrt{\Delta_{\rm fold}}}{4a}.
\]
Then the pair \((u,v)=(\xi_+,\xi_-)\) satisfies the two symmetric
relations in \eqref{eq:fold-characteristic}:
\[
 \xi_++\xi_-=\frac{C_0}{2a},\qquad
 \mathfrak b=a^2(4\xi_+\xi_-+2).
\]
We use the root presentation on these real sheets and the polynomial
\(\Pi\) itself for the global transfer and generating identities.
We will obtain the following formulas from a finite-dimensional transfer
recurrence.  For a differentiable real function \(f\),
write
\[
 [f]_{u,v}:=
 \begin{cases}
  \dfrac{f(u)-f(v)}{u-v},&u\neq v,\\[2mm]
  f'(u),&u=v
 \end{cases}.
\]
Whenever the real pair \(u,v\) satisfies the two symmetric relations in
\eqref{eq:fold-characteristic}, the folded determinants are
\begin{equation}\label{eq:fold-even}
 q_{2m}(x,s)=a^m[\mathfrak f_m]_{u,v},\qquad
 \mathfrak f_m(t)=(t-\zeta_s)(U_m(t)+U_{m-1}(t)),
\end{equation}
and
\begin{equation}\label{eq:fold-odd}
 q_{2m+1}(x,s)=a^m[\mathfrak g_m]_{u,v},\qquad
 \mathfrak g_m(t)=\left((x-s)t-
 \frac{x(1+a^2)+2as}{2a}\right)U_m(t).
\end{equation}
The recurrence and generating identities are polynomial and hence global.
The real-pair formulas apply when the two symmetric relations hold, with the
displayed derivative as the collision value at \(u=v\).  When
\(\Delta_{\rm fold}<0\), the polynomial identities avoid choosing a complex
square-root branch.

\subsubsection{Derivation of the folded formulas}

Pair the coordinates in the order
\((1,n),(2,n-1),\ldots\).  In this order the symmetric matrix
\((xI_n-A_n)J_n-sI_n\) is block tridiagonal.  Its interior diagonal and
off-diagonal blocks are
\[
 D_{\rm int}=\begin{bmatrix}-s&x\\x&-s\end{bmatrix},
 \qquad
 E_{\rm off}=\begin{bmatrix}0&-1\\-a&0\end{bmatrix},
\]
with lower block \(E_{\rm off}^T\).  For even order the terminal block is
\[
 D_{\rm term}^{\rm e}=\begin{bmatrix}-1-s&x\\x&-a-s\end{bmatrix};
\]
for odd order the terminal block is the scalar \(x-s\), coupled to the
preceding block by the column \((-1,-a)^T\).  Since
\[
 (xI_n-A_n-sJ_n)=\bigl((xI_n-A_n)J_n-sI_n\bigr)J_n,
\]
the determinant of this block matrix, multiplied by \(\det J_n\), is
\(q_n(x,s)\).  Set \(\mathcal Q_m^{\rm e}=q_{2m}(x,s)\) and
\(\mathcal Q_m^{\rm o}=q_{2m+1}(x,s)\).  We derive a recurrence by keeping track of five minors.
For \(\tau\in\{\mathrm e,\mathrm o\}\), let \(K_m^\tau\) denote the
corresponding matrix
\[
 (xI_n-A_n-sJ_n)J_n=(xI_n-A_n)J_n-sI_n,
 \qquad n=2m\ \text{or}\ 2m+1,
\]
after simultaneously permuting its rows and columns into the paired
order.  Thus \(K_m^{\mathrm e}\) has order \(2m\),
\(K_m^{\mathrm o}\) has order \(2m+1\), and
\[
 K_1^{\mathrm e}=
 \begin{bmatrix}-1-s&x\\x&-a-s\end{bmatrix},
 \qquad
 K_1^{\mathrm o}=
 \begin{bmatrix}
 -s&x&-1\\
 x&-s&-a\\
 -1&-a&x-s
 \end{bmatrix}.
\]
For \(m\geq1\), increasing \(m\) by one adds one interior block at the
exposed end:
\begin{equation}\label{eq:fold-block-extension}
 K_{m+1}^\tau=
 \begin{bmatrix}
  D_{\rm int}&\mathcal E_m^\tau\\
  (\mathcal E_m^\tau)^T&K_m^\tau
 \end{bmatrix},
 \qquad
 \mathcal E_m^\tau=
 \begin{bmatrix}E_{\rm off}&0\end{bmatrix},
\end{equation}
where the zero block has the width required by \(K_m^\tau\).

For a matrix whose first two coordinates are exposed, write
\(K(\widehat I\mid\widehat J)\) for the matrix obtained by deleting
rows \(I\) and columns \(J\).  Suppressing the parity superscript, set
\[
\begin{aligned}
 \mathfrak d_m&=\det K_m,&
 \mathfrak p_m&=\det K_m(\widehat1\mid\widehat1),&
 \mathfrak q_m&=\det K_m(\widehat2\mid\widehat2),\\
 \mathfrak c_m&=\det K_m(\widehat1\mid\widehat2),&
 \mathfrak r_m&=\det K_m(\widehat{1,2}\mid\widehat{1,2}).
\end{aligned}
\]
The off-principal minor obtained by deleting row \(2\) and column \(1\)
has the same value \(\mathfrak c_m\), because \(K_m\) is symmetric.

The one-step transfer follows from a Schur complement.  On the dense set
where both \(K_m\) and the tail obtained by deleting its first two
coordinates are invertible, let
\[
 F=\begin{bmatrix}u&w\\w&v\end{bmatrix}
\]
be its Schur complement in \(K_m\), and let \(\delta\) be the
determinant of that tail.  Then
\[
 (\mathfrak d_m,\mathfrak p_m,\mathfrak q_m,\mathfrak c_m,\mathfrak r_m)
 =\delta(uv-w^2,v,u,w,1).
\]
Since
\[
 E_{\rm off}\operatorname{adj}(F)E_{\rm off}^T
 =
 \begin{bmatrix}u&-aw\\-aw&a^2v\end{bmatrix},
\]
taking the Schur complement of \(K_m\) in
\eqref{eq:fold-block-extension} gives
\begin{equation}\label{eq:five-minor-update}
\begin{aligned}
 \mathfrak d_{m+1}
 &= (s^2-x^2)\mathfrak d_m
    +s(a^2\mathfrak p_m+\mathfrak q_m)
    -2ax\mathfrak c_m+a^2\mathfrak r_m,\\
 \mathfrak p_{m+1}&=-s\mathfrak d_m-a^2\mathfrak p_m,\\
 \mathfrak q_{m+1}&=-s\mathfrak d_m-\mathfrak q_m,\\
 \mathfrak c_{m+1}&=x\mathfrak d_m+a\mathfrak c_m,\\
 \mathfrak r_{m+1}&=\mathfrak d_m.
\end{aligned}
\end{equation}
Both sides are polynomial in \(x\) and \(s\).  For each fixed \(0<a<1\),
the identities hold on a dense nonsingular set, so the polynomial identity
theorem extends them to all \(x,s\), including singular matrix cases.

For the two terminal blocks, direct evaluation gives
\[
\begin{aligned}
&(\mathfrak d_1,\mathfrak p_1,\mathfrak q_1,\mathfrak c_1,\mathfrak r_1)_{\mathrm e}\\
&\qquad=\bigl((1+s)(a+s)-x^2,-a-s,-1-s,x,1\bigr),\\
&(\mathfrak d_1,\mathfrak p_1,\mathfrak q_1,\mathfrak c_1,\mathfrak r_1)_{\mathrm o}\\
&\qquad=\bigl(x(1+a^2)+2as-C_0(x-s),\\
&\hspace{26mm}s^2-sx-a^2,\ s^2-sx-1,\ x^2-sx-a,\ x-s\bigr).
\end{aligned}
\]
Extend the five minors to \(m=0\) by the formal initial values
\[
 (\mathfrak d_0,\mathfrak p_0,\mathfrak q_0,\mathfrak c_0,\mathfrak r_0)_{\mathrm e}
 =(1,a^{-1},1,0,a^{-1}),
 \qquad
 (\mathfrak d_0,\mathfrak p_0,\mathfrak q_0,\mathfrak c_0,\mathfrak r_0)_{\mathrm o}
 =(x-s,1,1,-1,0).
\]
One step of \eqref{eq:five-minor-update} reproduces the terminal minors.
For \(v\in\{d,p,q,c,r\}\), let
\[
 \mathcal F_v(z)=\sum_{m\geq0}(-1)^m\mathfrak v_mz^m.
\]
Since \(\det J_{2m}=\det J_{2m+1}=(-1)^m\),
\(\mathcal F_d(z)=\sum_{m\geq0}\mathcal Q_m^\tau z^m\).
Summing \eqref{eq:five-minor-update} with these signs gives
\[
 \mathcal F_p=\frac{\mathfrak p_0+sz\mathcal F_d}{1-a^2z},\qquad
 \mathcal F_q=\frac{\mathfrak q_0+sz\mathcal F_d}{1-z},\qquad
 \mathcal F_c=\frac{\mathfrak c_0-xz\mathcal F_d}{1+az},\qquad
 \mathcal F_r=\mathfrak r_0-z\mathcal F_d,
\]
and
\[
 [1-(x^2-s^2)z-a^2z^2]\mathcal F_d
 +sz(a^2\mathcal F_p+\mathcal F_q)-2axz\mathcal F_c
 =\mathfrak d_0-a^2z\mathfrak r_0.
\]
Substitute the first three expressions into the last identity and multiply
by \((1-a^2z)(1-z)(1+az)\).  Collecting terms gives
\[
 (1-az)\Pi(z)\mathcal F_d(z)=(1-az)\mathcal N^\tau(z),
\]
where the two initial states give
\[
 \mathcal N^{\mathrm e}(z)=(1+az)(1-2a\zeta_sz+a^2z^2),\qquad
 \mathcal N^{\mathrm o}(z)=(x-s)-[x(1+a^2)+2as]z+a^2(x-s)z^2.
\]
All denominators have constant term one and are therefore invertible in
\(\mathbb R[[z]]\).  Cancelling \(1-az\) yields
\begin{equation}\label{eq:fold-generating-even}
 \sum_{m\geq0}\mathcal Q_m^{\rm e}z^m
 =(1+az)(1-2a\zeta_sz+a^2z^2)\Pi(z)^{-1},
\end{equation}
and
\begin{equation}\label{eq:fold-generating-odd}
 \sum_{m\geq0}\mathcal Q_m^{\rm o}z^m
 =\bigl((x-s)-[x(1+a^2)+2as]z+a^2(x-s)z^2\bigr)
  \Pi(z)^{-1}.
\end{equation}

Coefficient extraction converts these generating functions to the
divided-difference formulas used in the half-angle chord parametrization.
Fix \(0<a<1\) and set \(U_{-1}=0\).  The Chebyshev generating function gives
\[
 \sum_{m\geq0}(u-\zeta_s)(U_m(u)+U_{m-1}(u))\omega_{\rm gf}^m
 =\frac{(u-\zeta_s)(1+\omega_{\rm gf})}
 {1-2u\omega_{\rm gf}+\omega_{\rm gf}^2},
\]
and, with \(\mathsf B_s^{\rm fold}=[x(1+a^2)+2as]/(2a)\),
\[
 \sum_{m\geq0}((x-s)u-\mathsf B_s^{\rm fold})U_m(u)\omega_{\rm gf}^m
 =\frac{(x-s)u-\mathsf B_s^{\rm fold}}
 {1-2u\omega_{\rm gf}+\omega_{\rm gf}^2}.
\]
Let \(u,v\in\mathbb R\) satisfy the two symmetric relations in
\eqref{eq:fold-characteristic}.  Taking the corresponding divided
difference of the Chebyshev generating identities and putting
\(\omega_{\rm gf}=az\) recovers
\eqref{eq:fold-generating-even} and \eqref{eq:fold-generating-odd}.
Coefficient comparison gives \eqref{eq:fold-even}--\eqref{eq:fold-odd},
including the derivative value when \(u=v\).  The underlying polynomial
identities remain valid when \(\Delta_{\rm fold}<0\), without choosing a
square-root branch.


\subsection{Half-angle chord parametrization}
\label{subsec:half-angle-chords}
Put \(h=-\tfrac12\log a>0\), so that \(a=e^{-2h}\), and set
\[
 \mathfrak s_h=\sinh h,\qquad \mathfrak c_h=\cosh h.
\]
For \(K\geq2\), define
\begin{equation}\label{eq:Phi-Psi-def}
\begin{aligned}
 \Phi_K(\theta)
 &=\sqrt{\mathfrak s_h^2+\sin^2\theta}\,
   \frac{|\sin(K\theta)|}{\sin\theta},
   &&0\leq\theta\leq\pi,\\
 \Psi_K(\eta)
 &=\sqrt{\mathfrak s_h^2-\sinh^2\eta}\,
   \frac{\sinh(K\eta)}{\sinh\eta},
   &&0\leq\eta\leq h.
\end{aligned}
\end{equation}
The quotients in \eqref{eq:Phi-Psi-def} are interpreted continuously at
\(\theta=0,\pi\) and \(\eta=0\).  The following half-angle reduction also
fixes the sign convention.  Put
\[
 y=\cos\theta,\qquad z=\cosh\eta,\qquad
 \mathfrak p_\theta=\sqrt{\mathfrak c_h^2-y^2},\qquad
 \mathfrak p_\eta=\sqrt{\mathfrak c_h^2-z^2},
\]
and let \(L_0=2r/\mathfrak c_h\).  If
\(x=L_0yz\), \(s=L_0\mathfrak p_\theta\mathfrak p_\eta>0\), and
\(\tau=\varsigma s\) with \(\varsigma\in\{-1,1\}\), then direct
expansion gives
\begin{equation}\label{eq:half-variable-identities}
\begin{gathered}
 (1+a+x)^2-s^2=4a(z+y)^2,\qquad
 (1+a-x)^2-s^2=4a(z-y)^2,\\
 u=2z^2-1,\qquad v=2y^2-1,\qquad
 \zeta_\tau=2\mathfrak c_h^2-1
 +2\varsigma\mathfrak p_\theta\mathfrak p_\eta.
\end{gathered}
\end{equation}
Substitution of \eqref{eq:half-variable-identities} shows that the real pair
\((u,v)\) satisfies the two symmetric relations in
\eqref{eq:fold-characteristic}.
For \(K=2m+1\), use
\[
 U_m(2t^2-1)+U_{m-1}(2t^2-1)=U_{K-1}(t);
\]
for \(K=2m+2\), use
\[
 U_m(2t^2-1)=\frac{U_{K-1}(t)}{2t}.
\]
Substitution of this pair in \eqref{eq:fold-even} and
\eqref{eq:fold-odd}, respectively, reduces either equation
\(q_{K-1}(x,\tau)=0\) to
\begin{equation}\label{eq:signed-half-chord}
 \mathfrak p_\theta U_{K-1}(y)
 =\varsigma\mathfrak p_\eta U_{K-1}(z).
\end{equation}
In the even-fold case the intermediate factorization is
\[
 (\mathfrak p_\eta+\varsigma\mathfrak p_\theta)
 \bigl(\mathfrak p_\eta U_{K-1}(z)
 -\varsigma\mathfrak p_\theta U_{K-1}(y)\bigr)=0.
\]
The first factor could vanish only if \(\varsigma=-1\) and
\(\mathfrak p_\eta=\mathfrak p_\theta\), which would imply
\(y^2=z^2\); the distinct-variable condition therefore leaves precisely
\eqref{eq:signed-half-chord}.  In the odd-fold case the reduction also
contains the factor
\(\mathfrak p_\eta y+\varsigma\mathfrak p_\theta z\).  On a
distinct-variable sheet this factor is nonzero: if it vanished, squaring
would give
\[
 (\mathfrak c_h^2-z^2)y^2=(\mathfrak c_h^2-y^2)z^2,
\]
and hence \(y^2=z^2\).  Thus the factor is nonzero whenever the folded
variables are distinct.  On the outer lobe
\(U_{K-1}(z)>0\), whereas, for
\(j\pi/K<\theta<(j+1)\pi/K\),
\[
 \operatorname{sgn}\bigl(U_{K-1}(\cos\theta)\bigr)
 =\operatorname{sgn}\bigl(\sin(K\theta)\bigr)=(-1)^j.
\]
Moreover,
\(\mathfrak p_\theta|U_{K-1}(y)|=\Phi_K(\theta)\) and
\(\mathfrak p_\eta U_{K-1}(z)=\Psi_K(\eta)\).  Thus
\eqref{eq:signed-half-chord} is equivalent to
\(\varsigma=(-1)^j\) and
\begin{equation}\label{eq:chord-equation}
 \Phi_K(\theta)=\Psi_K(\eta),
\end{equation}
with
\begin{equation}\label{eq:channel-xs}
 s=\mathcal S_{\rm hyp}(\theta,\eta)
 :=\frac{2r}{\mathfrak c_h}
 \sqrt{(\mathfrak s_h^2+\sin^2\theta)
       (\mathfrak s_h^2-\sinh^2\eta)},\qquad
 x=\frac{2r}{\mathfrak c_h}\cos\theta\cosh\eta.
\end{equation}
Consequently \(q_{K-1}(x,(-1)^j s)=0\).  Conversely, for distinct real
folded variables, their sum and product determine
\[
 x^2=L_0^2y^2z^2,\qquad
 s^2=L_0^2(\mathfrak c_h^2-y^2)(\mathfrak c_h^2-z^2).
\]
Thus, for distinct real variables and a fixed sign of \(x\), the two chord
relations are equivalent to the corresponding signed equation.  The
continuous divided-difference convention extends the determinant identities
to a folded-variable collision; the selected middle branch is shown below to
be collision-free.
Here \(K-1\) is the matrix size.  If the exterior argument has crossed
through \(1\), write it as \(z=\cos\phi\),
\(0\leq\phi<\pi/K\).  Then the chord equation is
\(\Phi_K(\theta)=\Phi_K(\phi)\), and the second factor in the formula
for \(s\) is \(\mathfrak s_h^2+\sin^2\phi\).  We denote the resulting
value by
\[
 \mathcal S_{\rm ell}(\theta,\phi)
 =\frac{2r}{\mathfrak c_h}
 \sqrt{(\mathfrak s_h^2+\sin^2\theta)
       (\mathfrak s_h^2+\sin^2\phi)}.
\]

\subsection{Outer-lobe geometry and chord classification}
\label{subsec:outer-lobe}
To identify the chords that realize the middle branch of
\cref{lem:middle-branch}, set
\begin{equation}\label{eq:Wcal-def}
 \mathcal W_K(u)=\sqrt{\mathfrak c_h^2-u^2}\,|U_{K-1}(u)|,
 \qquad |u|\leq\mathfrak c_h.
\end{equation}
At \(u=\pm\mathfrak c_h\), \eqref{eq:Wcal-def} denotes the continuous
endpoint value \(0\).
Then
\[
 \Phi_K(\theta)=\mathcal W_K(\cos\theta),\qquad
 \Psi_K(\eta)=\mathcal W_K(\cosh\eta).
\]
If \(u_{\ell,K}^{(0)}=\cos\frac{\ell\pi}{K}\), then, away from the zeros,
\[
 \frac{d}{du}\log\mathcal W_K(u)
 =\sum_{\ell=1}^{K-1}\frac1{u-u_{\ell,K}^{(0)}}
 -\frac{u}{\mathfrak c_h^2-u^2},
\]
and its derivative is
\begin{equation}\label{eq:W-log-concave}
 -\sum_{\ell=1}^{K-1}\frac1{(u-u_{\ell,K}^{(0)})^2}
 -\frac{\mathfrak c_h^2+u^2}{(\mathfrak c_h^2-u^2)^2}<0.
\end{equation}
The strict negativity in \eqref{eq:W-log-concave} makes the logarithmic
derivative strictly decreasing.  At each finite nodal endpoint it tends to the
appropriate signed infinity, and it tends to \(-\infty\) as
\(u\uparrow\mathfrak c_h\).  Thus there is exactly one lobe maximum in
every nodal
interval, including the outer interval
\(\mathcal J_K^{\rm out}=(\cos\frac{\pi}{K},\mathfrak c_h)\).  Let
\(u_K^{\max}\) be the unique maximizer on \(\mathcal J_K^{\rm out}\).
For every
\(0\leq\ell<\mathcal W_K(u_K^{\max})\), there is a unique solution
\[
 z_K(\ell)\in(u_K^{\max},\mathfrak c_h],\qquad
 \mathcal W_K(z_K(\ell))=\ell;
\]
we call it the \emph{endpoint-side outer solution}.  If
\(z_K(\ell)\geq1\), write it as \(\cosh\eta\),
\(0\leq\eta\leq h\); if it is at most \(1\), write it as
\(\cos\phi\), \(0\leq\phi<\pi/K\).  These are two coordinates on the
same outer lobe, joined at \(z=1\).

For a positive-half gap and
\(\theta\in(j\pi/K,(j+1)\pi/K)\), \(\theta\leq\pi/2\), a
\emph{$j$-chord} is the pair
\[
 y=\cos\theta,\qquad
 z=z_K(\mathcal W_K(y)),\qquad
 \mathcal W_K(y)=\mathcal W_K(z),
\]
together with the coordinate-free values
\[
 x=L_0yz,\qquad
 s=L_0\sqrt{(\mathfrak c_h^2-y^2)(\mathfrak c_h^2-z^2)};
\]
its signed root is
\((-1)^j s\).  The term endpoint-side always refers to the part of the
outer lobe adjacent to \(\mathfrak c_h\), independently of whether it is
described by the hyperbolic or elliptic coordinate.

Reflection reduces the branch classification to a positive-half gap, so
\(j\geq1\)
and \(\theta\leq\pi/2\).  Put
\(\phi_{\rm rem}=\theta-j\pi/K\in(0,\pi/K)\).  Since
\(|\sin(K\phi_{\rm rem})|=|\sin(K\theta)|\) and
\[
 \mathcal A_h(t)=
 \frac{\sqrt{\mathfrak s_h^2+\sin^2t}}{\sin t}
\]
is strictly decreasing on \((0,\pi/2]\),
\begin{equation}\label{eq:outer-level-dominates}
 \Phi_K(\theta)<\Phi_K(\phi_{\rm rem})
 \leq\max_{u\in\mathcal J_K^{\rm out}}\mathcal W_K(u).
\end{equation}
Thus the endpoint-side outer solution exists for every inner parameter.

At \(x=\lambda_{K-1,k}=2r\cos\frac{k\pi}{K}\), \(s=0\), the two
folded variables are \(\cos\frac{2k\pi}{K}\) and
\(\cosh(2h)\), or, in half variables, \(\cos\frac{k\pi}{K}\) and
\(\mathfrak c_h\).  Moreover, zero is a simple eigenvalue of
\(B_{K-1}(\lambda_{K-1,k})\): the matrix
\(\lambda_{K-1,k}I_{K-1}-A_{K-1}\) has rank \(K-2\), and \(J_{K-1}\) is
invertible.  The implicit-function theorem therefore gives, near either
gap endpoint, a real continuation of the middle signed root and real
half variables \(y=\cos\theta\) and \(z\).

\subsection{Exhaustion of the selected chord sheet}
\label{subsec:chord-exhaustion}

\begin{lemma}[Endpoint-side chord-sheet exhaustion]
\label[lemma]{lem:chord-sheet-exhaustion}
Let \(K\geq3\) and let \(j\geq1\) be an integer.  Under the conventions of
\cref{eq:chord-equation,eq:channel-xs}, if
\[
 \frac{(j+1)\pi}{K}\leq\frac{\pi}{2},
\]
then every point of the ordered middle branch in the corresponding open
positive gap has a selected endpoint-side \(j\)-chord representation with
\(j\pi/K<\theta<(j+1)\pi/K\), and every \(\theta\) in this interval gives a
point of that branch.  The signed root is \((-1)^j s\), its magnitude is
\(s_{\min}(xI_{K-1}-A_{K-1})\), and the two folded variables are distinct.
This includes a gap adjacent to zero.  Separately, if \(K\) is odd and
\(j=(K-1)/2\), the same bidirectional statement holds from the positive
spectral endpoint to the center of the single central gap, with
\(j\pi/K<\theta\leq\pi/2\).
\end{lemma}

\begin{proof}
Continue the ordered middle root \(s_*(x)\) through the open gap.  Write
\(\Delta_*(x)=\Delta_{\rm fold}(x,s_*(x))\) for its folded discriminant,
let \(y(x)\leq z(x)\) be the reconstructed nonnegative half variables on
the real sheet issuing from the endpoint germ, and set
\(\theta_{\rm ch}(x)=\arccos y(x)\).  Let \(\mathcal C\) be the set of gap points for
which
\[
 \Delta_*(x)>0,\qquad
 \frac{j\pi}{K}<\theta_{\rm ch}(x)<\frac{(j+1)\pi}{K},\qquad
 u_K^{\max}<z(x)<\mathfrak c_h.
\]
On \(\mathcal C\), the reversible folded determinant identities give
\(\mathcal W_K(y(x))=\mathcal W_K(z(x))\); uniqueness on the endpoint-side
outer lobe then gives
\[
 z(x)=z_K(\mathcal W_K(y(x))).
\]
Thus \(\mathcal C\) is precisely the selected endpoint-side chord sheet.
At the upper gap endpoint, \(z=\mathfrak c_h\) and
\(\theta_{\rm ch}=j\pi/K\).  There
\(z=\mathfrak c_h>1\geq y\), so \(\Delta_*>0\) persists immediately
inside the gap.  Moreover,
\cref{lem:middle-branch} gives \(\operatorname{sgn}s_*=(-1)^j\); the signed
chord equation therefore places
\(\theta_{\rm ch}\) in
\((j\pi/K,(j+1)\pi/K)\).  Since \(|s_*|>0\), the channel identity gives
\(z<\mathfrak c_h\), while continuity keeps \(z>u_K^{\max}\).  Hence
\(\mathcal C\) is nonempty.  The ordered eigenvalues of the continuous real symmetric pencil
\(B_{K-1}(x)\) are continuous in \(x\), including at an eigenvalue crossing,
so the continuation passes through a multiple least singular value.  The
folded pair \(u=2z^2-1\), \(v=2y^2-1\) and its reconstructed half variables
are continuous, while the sheet inequalities are strict.  Thus
\(\mathcal C\) is relatively open in the gap.

It is also relatively closed.  At a relative closure point the chord
identity and all sheet inequalities persist weakly.  The possible closure
cases are as follows.  (1) If \(\Delta_*(x)=0\), then the reconstructed half
variables coincide, which is impossible because
\[
 z\geq u_K^{\max}>\cos\frac{\pi}{K}
 \geq\cos\frac{j\pi}{K}\geq y.
\]
(2) At an interior point of a positive gap, including one adjacent to zero,
both half variables retain their signs: \(z>0\), while \(y=0\) would force
\(x=0\), outside the open positive half.  (3) Neither the outer variable can
reach \(\mathfrak c_h\) nor the inner variable a nodal endpoint: in either
case the common level is zero, the endpoint-side solution is
\(z=\mathfrak c_h\), and \eqref{eq:channel-xs} gives \(s=0\), impossible
inside a spectral gap.  (4) The
outer variable remains in its endpoint-side component, since its first exit
point would be \(u_K^{\max}\), contradicting
\eqref{eq:outer-level-dominates}.  Passage through \(z=1\) changes only the
coordinate representation.  Hence every relative closure point remains in
\(\mathcal C\).  Since the open gap is an interval, \(\mathcal C\) is the
entire gap.  This proves the first exhaustion direction: every actual gap
point lies on the selected endpoint-side chord sheet.

For the converse direction, the reconstructed angle \(\theta_{\rm ch}\) extends
continuously to the closed spectral gap.  The exact folded variables at the
two endpoints give
\[
 \theta_{\rm ch}(\lambda_{K-1,j+1})=\frac{(j+1)\pi}{K},\qquad
 \theta_{\rm ch}(\lambda_{K-1,j})=\frac{j\pi}{K}.
\]
For every \(\theta_0\in(j\pi/K,(j+1)\pi/K)\), the intermediate value
theorem therefore supplies an interior gap point \(x\) with
\(\theta_{\rm ch}(x)=\theta_0\).  The preceding classification and uniqueness of
the endpoint-side outer solution then give
\[
 y(x)=\cos\theta_0,\qquad
 z(x)=z_K(\mathcal W_K(\cos\theta_0)),
\]
so the corresponding abscissa and signed root are exactly those of the
prescribed \(j\)-chord.  Thus every chord parameter produces an actual point
of the ordered middle branch.  By \cref{lem:middle-branch}, its positive
magnitude is \(s_{\min}(xI_{K-1}-A_{K-1})\).

When \(K\) is odd, the same argument applies to the positive half of
the single central gap, continuing from the positive spectral endpoint
to \(x=0\).  Because \(z>0\), the identity
\(x=(2r/\mathfrak c_h)\cos\theta\,z\) forces \(\theta=\pi/2\) at
\(x=0\).  Thus every parameter on that half occurs, and reflection covers
the negative half.
\end{proof}


\subsection{Noncentral gap comparison}
\label{subsec:noncentral-comparison}

\begin{lemma}[Noncentral gap comparison]
\label[lemma]{lem:noncentral-gap-comparison}
Let \(j\geq1\) and \(N\geq2(j+1)\) be integers, and put \(M=N+1\).
For every \(x\in G_{M-1,j}^{\circ}\), there is
\(\widehat x\in G_{N-1,j}^{\circ}\) such that
\[
 g_{N-1}(\widehat x)>g_{M-1}(x).
\]
\end{lemma}

\begin{proof}
Fix \(x\in G_{M-1,j}^{\circ}\).  By
\cref{lem:chord-sheet-exhaustion}, its middle branch has a hyperbolic or
elliptic chord representation with
\[
 \theta\in\left(\frac{j\pi}{M},\frac{(j+1)\pi}{M}\right).
\]
Set \(\kappa=N/M\) and \(\theta_0=\theta/\kappa\).  The hypotheses place
\(\theta_0\) in the corresponding \(N\)-nodal interval, with every
trigonometric argument below in \((0,\pi/2]\).

First suppose \(\Phi_M(\theta)=\Psi_M(\eta)\).  Concavity of sine gives
\begin{equation}\label{eq:Phi-scale}
 \frac{\Phi_N(\theta_0)}{\Phi_M(\theta)}
 =\frac{\mathcal A_h(\theta_0)}{\mathcal A_h(\theta)}>\kappa.
\end{equation}
Indeed, \(\theta=\kappa\theta_0\) and strict concavity gives
\(\sin\theta>\kappa\sin\theta_0\); hence
\[
 \mathcal A_h(\theta_0)^2-\kappa^2\mathcal A_h(\theta)^2
 =(1-\kappa^2)+\mathfrak s_h^2\left(
 \frac1{\sin^2\theta_0}-\frac{\kappa^2}{\sin^2\theta}
 \right)>0.
\]
On the other hand,
\begin{equation}\label{eq:Psi-scale}
 \frac{\Psi_N(\eta)}{\Psi_M(\eta)}
 =\frac{\sinh(N\eta)}{\sinh(M\eta)}\leq\kappa,
\end{equation}
with strict inequality for \(\eta>0\), by convexity of \(\sinh\) on
the positive half-axis.  Together, \eqref{eq:Phi-scale} and
\eqref{eq:Psi-scale} imply
\(\Phi_N(\theta_0)>\Psi_N(\eta)\).  Since \(\Phi_N\) vanishes at the right
endpoint \((j+1)\pi/N\), its rightmost solution
\(\widehat\theta\) of
\[
 \Phi_N(\widehat\theta)=\Psi_N(\eta)
\]
exists by the intermediate value theorem and satisfies
\(\widehat\theta>\theta_0>\theta\).  Strict log-concavity of
\(\mathcal W_N\) makes it the unique solution on the descending half
of that lobe.  The same exterior
variable is endpoint-side for the \(N\)-problem as well.  Indeed, for
\(u=\cosh\eta>1\),
\[
 \frac{d}{du}\log\mathcal W_M(u)-
 \frac{d}{du}\log\mathcal W_N(u)
 =\frac{M\coth(M\eta)-N\coth(N\eta)}{\sinh\eta}>0.
\]
The numerator is positive because, for \(\eta>0\),
\[
 \frac{d}{dt}\bigl(t\coth(t\eta)\bigr)
 =\frac{\sinh(t\eta)\cosh(t\eta)-t\eta}{\sinh^2(t\eta)}>0.
\]
The endpoint-side branch is characterized by a nonpositive logarithmic
derivative.  The displayed difference therefore makes the
\(N\)-derivative strictly negative.  The case \(\eta=0\) follows by
continuity.
The chord classification above says that the constructed value
\(\mathcal S_{\rm hyp}(\widehat\theta,\eta)\) is the least singular
value at its constructed abscissa.  Formula
\eqref{eq:channel-xs} and the strict increase of sine on the interval
now give
\begin{equation}\label{eq:hyper-height-increase}
 \mathcal S_{\rm hyp}(\widehat\theta,\eta)
 >\mathcal S_{\rm hyp}(\theta,\eta).
\end{equation}

Suppose next that the exterior endpoint is elliptic, so that
\(\Phi_M(\theta)=\Phi_M(\phi)\), \(0<\phi<\pi/M\).  (The case
\(\phi=0\)
is the limiting hyperbolic case already covered.)  Put
\[
 \theta_0=\theta/\kappa,\qquad \phi_0=\phi/\kappa.
\]
The ratio
\[
 \mathcal A_{N,M}^{\Phi}(t)
 =\frac{\Phi_N(t/\kappa)}{\Phi_M(t)}
 =\frac{\mathcal A_h(t/\kappa)}{\mathcal A_h(t)}
\]
is strictly increasing while \(0<t<t/\kappa\leq\pi/2\).  Indeed,
\[
 \frac{d}{dt}\log\mathcal A_{N,M}^{\Phi}(t)
 =\mathfrak d_h(t)-\kappa^{-1}\mathfrak d_h(t/\kappa)>0,\qquad
 \mathfrak d_h(t)=\frac{\mathfrak s_h^2\cot t}
 {\mathfrak s_h^2+\sin^2t},
\]
because
\[
 t\mathfrak d_h(t)=(t\cot t)
 \frac{\mathfrak s_h^2}{\mathfrak s_h^2+\sin^2t}
\]
is strictly decreasing: both positive factors decrease on
\((0,\pi/2)\), and the second decreases strictly.  Since
\(\theta>\pi/M>\phi\), it follows that
\(\Phi_N(\theta_0)>\Phi_N(\phi_0)\).  The rightmost solution of
\[
 \Phi_N(\widehat\theta)=\Phi_N(\phi_0)
\]
therefore has \(\widehat\theta>\theta_0\).  Its exterior variable is
endpoint-side.  Indeed, if
\(\Lambda_K(t)=(\log\Phi_K)'(t)\), then on the
endpoint-side part of the first outer lobe
\(\Lambda_M(\phi)\geq0\), and differentiation of
\(\Phi_N(t/\kappa)=\mathcal A_{N,M}^{\Phi}(t)\Phi_M(t)\) gives
\[
 \kappa^{-1}\Lambda_N(t/\kappa)
 =\Lambda_M(t)
 +\frac{(\mathcal A_{N,M}^{\Phi})'(t)}{\mathcal A_{N,M}^{\Phi}(t)}>0.
\]
Because
\(\Lambda_K(\phi)=-\sin\phi\,(\log\mathcal W_K)'(\cos\phi)\),
positivity of \(\Lambda_N(\phi_0)\) is exactly the endpoint-side
condition for the
\(N\)-outer solution.  The one-maximum property now fixes the label, and the chord
classification says again that
\(\mathcal S_{\rm ell}(\widehat\theta,\phi_0)\) is the least singular
value at the constructed abscissa.  Here
\(\widehat\theta>\theta_0>\theta\) and \(\phi_0>\phi\), with all four
angles below \(\pi/2\).  Thus both sine factors in the elliptic version
of \eqref{eq:channel-xs} have increased, at least one strictly, and hence
\begin{equation}\label{eq:elliptic-height-increase}
 \mathcal S_{\rm ell}(\widehat\theta,\phi_0)
 >\mathcal S_{\rm ell}(\theta,\phi).
\end{equation}
Equations \eqref{eq:hyper-height-increase} and
\eqref{eq:elliptic-height-increase}, together with the coordinate formula
for \(x\), place the constructed abscissa \(\widehat x\) in
\(G_{N-1,j}^{\circ}\).  Chord exhaustion identifies the original and
constructed heights with \(g_{M-1}(x)\) and
\(g_{N-1}(\widehat x)\), respectively, proving the claim.
\end{proof}

\subsection{Central gap comparison}
\label{subsec:central-comparison}

For \(m\geq1\), define the even-order central height by
\begin{equation}\label{eq:c-height-def}
 \mathfrak h_m^{\mathrm e}=s_{\min}(aI_m+V_m).
\end{equation}
By reflection, the two gaps adjacent to the zero eigenvalue of
\(A_{2m+1}\) have the same height; denote it by
\begin{equation}\label{eq:d-height-def}
 \mathfrak h_m^{\mathrm o}=\max_{0\leq x\leq2r\sin\frac{\pi}{2m+2}}
 s_{\min}(xI_{2m+1}-A_{2m+1}).
\end{equation}
The central heights defined in \eqref{eq:c-height-def} and
\eqref{eq:d-height-def} satisfy the following strict alternation.

\begin{lemma}[Central gap comparison]\label[lemma]{lem:central-gap-comparison}
For every \(m\geq1\),
\begin{equation}\label{eq:central-interlace}
 \mathfrak h_m^{\mathrm e}>\mathfrak h_m^{\mathrm o}>
 \mathfrak h_{m+1}^{\mathrm e}.
\end{equation}
\end{lemma}

\begin{proof}
Put
\[
 \zeta=\frac{a+a^{-1}}2.
\]
For every \(j\geq1\), there is a unique
\(\rho_j>\cos\frac{\pi}{j+1}\) such that
\begin{equation}\label{eq:rho-central}
 aU_j(\rho_j)=U_{j-1}(\rho_j).
\end{equation}
Indeed, the Christoffel--Darboux identity gives, above the largest zero
of \(U_j\),
\[
 \frac{d}{dt}\frac{U_{j-1}(t)}{U_j(t)}
 =-\frac{2\sum_{k=0}^{j-1}U_k(t)^2}{U_j(t)^2}<0.
\]
The ratio decreases from \(+\infty\) to \(0\), proving existence and
uniqueness.  A three-term determinant recurrence gives
\begin{equation}\label{eq:c-determinant}
 \det\left((aI_j+V_j)^*(aI_j+V_j)
 -(1+a^2-2at)I_j\right)
 =a^j\left(U_j(t)-a^{-1}U_{j-1}(t)\right).
\end{equation}
Thus the largest root \(\rho_j\) gives the least eigenvalue.  Together
with
\[
 U_j(t)^2+U_{j-1}(t)^2-2tU_j(t)U_{j-1}(t)=1,
\]
equations \eqref{eq:rho-central} and \eqref{eq:c-determinant} give
\begin{equation}\label{eq:c-rho-identities}
 (\mathfrak h_j^{\mathrm e})^2
 =1+a^2-2a\rho_j=\frac1{U_j(\rho_j)^2},\qquad
 \zeta-\rho_j=\frac1{2U_{j-1}(\rho_j)U_j(\rho_j)}.
\end{equation}

We first identify \(\mathfrak h_m^{\mathrm e}\) as the central-gap height of
\(A_{2m}\).  Let
\[
 P_{\rm oe}=[\mathbf e_1,\mathbf e_3,\ldots,\mathbf e_{2m-1},
 \mathbf e_2,\mathbf e_4,\ldots,\mathbf e_{2m}]
\]
be the permutation matrix that places odd coordinates first.  Then
\[
 P_{\rm oe}^T A_{2m}P_{\rm oe}=
 \begin{bmatrix}0&I_m+aV_m^T\\aI_m+V_m&0\end{bmatrix}.
\]
Moreover,
\[
 (I_m+aV_m^T)^*(I_m+aV_m^T)
 -J_m(aI_m+V_m)^*(aI_m+V_m)J_m
 =(1-a^2)\mathbf e_m\mathbf e_m^T\geq0.
\]
The singular values of the block off-diagonal matrix are those of its two
off-diagonal blocks.  The Gram inequality shows that
\(s_{\min}(I_m+aV_m^T)\geq s_{\min}(aI_m+V_m)=
\mathfrak h_m^{\mathrm e}\), and hence
\(s_{\min}(A_{2m})=\mathfrak h_m^{\mathrm e}\).  To maximize over the central
gap, put \(K=2m+1\) and use the half variables
\[
 y=\cos\theta,\qquad
 z=\cosh\eta\quad\text{or}\quad z=\cos\phi.
\]
On the positive half of the central middle branch, the preceding chord
classification gives
\begin{equation}\label{eq:even-central-half}
 \mathcal W_K(y)=\mathcal W_K(z),\qquad
 x=\frac{2r}{\mathfrak c_h}yz,\qquad
 s^2=\frac{4a}{\mathfrak c_h^2}
 (\mathfrak c_h^2-y^2)(\mathfrak c_h^2-z^2),
\end{equation}
where \(0\leq y<\cos\frac{m\pi}{K}\) is in the central lobe and
\(z\) is the
endpoint-side outer solution.  The function \(\mathcal W_K\) is even,
and strict log-concavity on the central lobe gives
\[
 \mathcal W_K(y)<\mathcal W_K(0)=\mathfrak c_h\qquad(y>0).
\]
Let
\[
 z_0=\sqrt{\frac{1+\rho_m}{2}}.
\]
The half-angle identity and \eqref{eq:rho-central} give
\[
 U_{2m}(z_0)=U_m(\rho_m)+U_{m-1}(\rho_m)
 =(1+a)U_m(\rho_m).
\]
Together with \eqref{eq:c-rho-identities}, this yields
\begin{equation}\label{eq:z0-central}
 \mathcal W_K(z_0)=\mathfrak c_h,\qquad
 4a(\mathfrak c_h^2-z_0^2)=2a(\zeta-\rho_m)
 =(\mathfrak h_m^{\mathrm e})^2.
\end{equation}
At \(x=0\), \cref{lem:middle-branch} and chord exhaustion identify the nearest
root with \(\mathfrak h_m^{\mathrm e}\), so \(z_0\) is the endpoint-side outer
solution of level \(\mathfrak c_h\).  If \(x\neq0\), the level in
\eqref{eq:even-central-half} is strictly below \(\mathfrak c_h\), and its
endpoint-side solution satisfies \(z>z_0\).  Equations
\eqref{eq:even-central-half}--\eqref{eq:z0-central} then give
\[
 s^2<4a(\mathfrak c_h^2-z_0^2)=(\mathfrak h_m^{\mathrm e})^2.
\]
Reflection gives the same strict bound on the other half of the gap.
Therefore the central-gap height is \(\mathfrak h_m^{\mathrm e}\), uniquely
attained at \(x=0\).

Set
\[
 \omega_m^{\rm e}=\cos\frac{\pi}{m+1},\qquad
 \mathcal V_m(t)=\sqrt{(\zeta-t)(1+t)}\,U_m(t).
\]
For the central specialization, take \(K=2m+2\) and set
\[
 t_{\rm in}=2y^2-1=\cos(2\theta),\qquad
 t_{\rm out}=2z^2-1.
\]
On the positive central gap,
\(-1<t_{\rm in}<-\omega_m^{\rm e}\), while the endpoint-side outer variable
satisfies
\(\omega_m^{\rm e}<t_{\rm out}<\zeta\), where
\(\zeta=2\mathfrak c_h^2-1\).  The identity
\(U_{2m+1}(t)=2tU_m(2t^2-1)\) gives
\[
 \begin{aligned}
 \mathcal W_{2m+2}(y)^2
  &= (\zeta-t_{\rm in})(1+t_{\rm in})U_m(t_{\rm in})^2,\\
 \mathcal W_{2m+2}(z)^2
  &= (\zeta-t_{\rm out})(1+t_{\rm out})U_m(t_{\rm out})^2.
 \end{aligned}
\]
Because \(\zeta+1=2\mathfrak c_h^2\), the formulas
\(x=(2r/\mathfrak c_h)yz\) and
\(s^2=(4a/\mathfrak c_h^2)
(\mathfrak c_h^2-y^2)(\mathfrak c_h^2-z^2)\) become the following chord
parametrization:
\begin{equation}\label{eq:central-chord}
\begin{gathered}
 -1<t_{\rm in}<-\omega_m^{\rm e},\qquad
 \omega_m^{\rm e}<t_{\rm out}<\zeta,\qquad
 \mathcal V_m(t_{\rm out})=|\mathcal V_m(t_{\rm in})|,\\
 x^2=\frac{2a(1+t_{\rm out})(1+t_{\rm in})}{\zeta+1},\qquad
 s^2=\frac{2a(\zeta-t_{\rm out})(\zeta-t_{\rm in})}{\zeta+1},
\end{gathered}
\end{equation}
where \(t_{\rm out}\) is the descending, or largest, solution.  To derive
\eqref{eq:central-chord} from \eqref{eq:chord-equation}, use
\[
 \mathfrak s_h^2+\sin^2\theta=\frac{\zeta-t_{\rm in}}{2},\qquad
 \mathfrak s_h^2-\sinh^2\eta=\frac{\zeta-t_{\rm out}}{2}.
\]
Continuation through \(t_{\rm out}=1\) uses the elliptic coordinate
\(z=\cos\phi\).

Let \(\rho=\rho_m\).  The function \(\log\mathcal V_m\) is strictly
concave on \((\omega_m^{\rm e},\zeta)\).  If
\((\log\mathcal V_m)'(\rho)\geq0\), then \(\rho\) lies at or before the
unique lobe maximum.  The descending root lies at or after that maximum, so
\(t_{\rm out}\geq\rho\).  Now suppose
\((\log\mathcal V_m)'(\rho)<0\).  Then \(\rho>a\): if
\(\rho\leq a<1\), then \(\rho<1\), hence \(\rho^2-1\neq0\), and the
following quotient is well defined:
\[
 \frac{U_m'(\rho)}{U_m(\rho)}
 =\frac{m\rho-(m+1)a}{\rho^2-1}
\]
would give
\[
\begin{aligned}
 (\log\mathcal V_m)'(\rho)
 &=\frac{(2m+2)a+1-(2m+1)\rho}{2(1-\rho^2)}
   -\frac1{2(\zeta-\rho)}\\
 &\geq\frac1{2(1-\rho)}-\frac1{2(\zeta-\rho)}\geq0,
\end{aligned}
\]
which contradicts the assumed negative derivative.  Hence \(\rho>a\).
Writing \(t_{\rm in}=-\cos\phi_{\rm c}\),
\(0<\phi_{\rm c}<\pi/(m+1)\), gives
\[
\begin{aligned}
 |\mathcal V_m(t_{\rm in})|^2
 &=\frac{(\zeta+\cos\phi_{\rm c})
          \sin^2((m+1)\phi_{\rm c})}{1+\cos\phi_{\rm c}}\\
 &<\frac{\zeta+\omega_m^{\rm e}}{1+\omega_m^{\rm e}}
 <\frac{1+a}{2a}<\frac{1+\rho}{2a}=\mathcal V_m(\rho)^2.
\end{aligned}
\]
Now \(\rho\) is on the descending side and
\(|\mathcal V_m(t_{\rm in})|<\mathcal V_m(\rho)\), so the
unique descending root again satisfies \(t_{\rm out}>\rho\).  In either case,
\(t_{\rm out}\geq\rho\), and \(t_{\rm in}>-1\) makes the following
inequality strict:
\[
 s^2<2a(\zeta-\rho)=(\mathfrak h_m^{\mathrm e})^2.
\]
This holds along the whole central sheet, so
\(\mathfrak h_m^{\mathrm o}<\mathfrak h_m^{\mathrm e}\).
Thus the odd-order central height lies strictly below the preceding
even-order central height.  We next compare it with the following even-order
height.

\proofstage{Lower comparison
\(\mathfrak h_m^{\mathrm o}>\mathfrak h_{m+1}^{\mathrm e}\): scalar estimate}
Put \(L=m+1\geq2\), and abbreviate
\[
 \rho=\rho_L,\quad u_-=U_{L-1}(\rho),\quad u_+=U_L(\rho),\quad
 \omega_L=\cos\frac{\pi}{2L},\quad \Delta_\omega=1-\omega_L,
\]
\begin{equation}\label{eq:central-D-xi}
 \mathfrak D_L=1+2(\rho+\omega_L)u_-u_+,\qquad
 \xi=\rho-\frac{\Delta_\omega}{\mathfrak D_L}.
\end{equation}
Thus \(a=u_-/u_+\) and \(\mathfrak h_L^{\mathrm e}=1/u_+\).

The required inequality is
\begin{equation}\label{eq:central-target}
 (1+\omega_L)(1+\xi)^2U_{L-1}(\xi)^2
 <(\rho+\omega_L)\mathfrak D_L.
\end{equation}

Put \(b=u_+/u_-=1/a\), and define
\[
 Q_b=b^2-2b\rho+1,\qquad S_b=b^2+2b\omega_L+1.
\]
The Chebyshev identity gives \(u_-^2Q_b=1\), and
\begin{equation}\label{eq:central-D-QS}
 \mathfrak D_L=S_b/Q_b,\qquad
 \Delta_\xi:=\rho-\xi=\Delta_\omega Q_b/S_b.
\end{equation}
Indeed, since \(u_+=bu_-\) and \(u_-^2=Q_b^{-1}\),
\[
 \mathfrak D_L=1+2(\rho+\omega_L)u_-u_+
  =1+\frac{2b(\rho+\omega_L)}{Q_b}=\frac{S_b}{Q_b},
\]
which proves the first identity in \eqref{eq:central-D-QS}; the second
follows from
\eqref{eq:central-D-xi}.
The ratio \(U_L/U_{L-1}\) is strictly increasing above the largest zero
of \(U_{L-1}\), and equals \(1\) at
\(\rho_{\rm ref}=\cos\frac{\pi}{2L+1}\).  Hence
\(\rho>\rho_{\rm ref}\).  Since
\(S_b-Q_b=2b(\rho+\omega_L)>0\),
\[
 \xi>\rho_{\rm ref}-\Delta_\omega>2\omega_L-1
 >2\omega_L^2-1=\cos\frac{\pi}{L}.
\]
Therefore \(f_L(x)=(1+x)U_{L-1}(x)\) is positive, strictly increasing,
and strictly log-concave on \([\xi,\rho]\).  In particular,
\(\mathfrak d_\rho=(\log f_L)'(\rho)>0\), and
\[
 (\rho^2-1)\mathfrak d_\rho=L(b-\rho)-1.
\]
This identity follows from the Chebyshev derivative recurrence and holds
also at \(\rho=1\).  Strict log-concavity gives
\[
 \log f_L(\rho)-\log f_L(\xi)>\Delta_\xi\mathfrak d_\rho,
 \qquad
 f_L(\xi)^2<f_L(\rho)^2e^{-2\Delta_\xi\mathfrak d_\rho}
 <\frac{f_L(\rho)^2}{1+2\Delta_\xi\mathfrak d_\rho}.
\]
Since
\(f_L(\rho)^2=(1+\rho)^2/Q_b\) and
\(\Delta_\xi=\Delta_\omega Q_b/S_b\), substitution into
\eqref{eq:central-target} reduces it to
\begin{equation}\label{eq:N-positive-target}
 \mathcal N=(\rho+\omega_L)
 (S_b+2\Delta_\omega Q_b\mathfrak d_\rho)
 -(1+\omega_L)(1+\rho)^2>0.
\end{equation}
Both cases use
\begin{equation}\label{eq:Delta-omega-basic-bound}
 \Delta_\omega=1-\cos\frac{\pi}{2L}
 <\frac{\pi^2}{8L^2}<\frac5{4L^2},
 \qquad \Delta_\omega L<1.
\end{equation}

If \(\rho\geq1\), put \(u=b-\rho\).  The derivative identity gives
\(u\geq1/L\).  Thus
\[
 \begin{aligned}
 S_b-(1+\rho)^2
 &=u(2\rho+u+2\omega_L)-2\rho\Delta_\omega\\
 &\geq2\rho(1/L-\Delta_\omega)+2\omega_L/L+1/L^2>0.
 \end{aligned}
\]
Together with \(\rho+\omega_L\geq1+\omega_L\) and
\(Q_b\mathfrak d_\rho>0\), this proves \eqref{eq:N-positive-target}.

If \(\rho<1\), set
\[
 \upsilon=1-\rho,\qquad \Delta_b=b-1,\qquad
 u=\Delta_b+\upsilon,\qquad \Delta_-=\upsilon(2-\upsilon).
\]
Since \(b>1\) and \(\mathfrak d_\rho>0\), the same derivative identity gives
\begin{equation}\label{eq:u-small-range}
 0<\upsilon<u<1/L.
\end{equation}
With \(\rho\) fixed, view \(\mathcal N\) as a function
\(\mathcal N_\rho^{<}(\Delta_b)\), substituting
\(\mathfrak d_\rho=(1-Lu)/\Delta_-\).
Expansion of \eqref{eq:N-positive-target} gives
\[
 \mathcal N_\rho^{<}(\Delta_b)-\mathcal N_\rho^{<}(0)
 =\frac{\Delta_b(\rho+\omega_L)}{\Delta_-}\mathfrak q(u),
\]
where
\[
 \mathfrak q(u)=\upsilon(2-\upsilon)
 [4-2\Delta_\omega+u-\upsilon]
 +2\Delta_\omega[(u+\upsilon)(1-Lu)-2L\upsilon].
\]
The bounds
\begin{equation}\label{eq:Delta-upsilon-bounds}
 \Delta_\omega<\frac5{4L^2},\qquad
 \upsilon<\frac5{(2L+1)^2}\leq\frac15
\end{equation}
follow from \eqref{eq:Delta-omega-basic-bound} and
\(\rho>\rho_{\rm ref}\).  Since \(L\geq2\), they imply
\[
 \mathfrak B_\rho:=(2-\Delta_\omega)(2-\upsilon)-4\Delta_\omega L
 >\frac{243}{80}-\frac52>0.
\]
Consequently,
\[
 \mathcal N_\rho^{<}(0)
 =\upsilon\mathfrak B_\rho
 +\frac{2\upsilon\Delta_\omega^2(2L-1)}{2-\upsilon}>0,
\]
and \eqref{eq:u-small-range} gives directly
\[
 \mathfrak q(u)
 >\upsilon[2(2-\Delta_\omega)(2-\upsilon)-4\Delta_\omega L]
 >\upsilon\mathfrak B_\rho>0.
\]
This proves \eqref{eq:N-positive-target}, and hence
\eqref{eq:central-target}, in both cases.

\proofstage{A point above the next even-order height}
Since \(a>0\), \(\Delta_\omega>0\), \(\rho+\omega_L>0\), and
\(\zeta+\omega_L>0\), let \(x_*>0\) be the positive square root determined
by
\begin{equation}\label{eq:x-star-central}
 x_*^2=\frac{2a\Delta_\omega(\rho+\omega_L)}{\zeta+\omega_L},
 \qquad \tau_*=\mathfrak h_L^{\mathrm e}.
\end{equation}
At \(x=x_*\) and \(s=\pm\tau_*\), where
\(\tau_*=\mathfrak h_L^{\mathrm e}\), the real pair
\(u=\xi\), \(v=-\omega_L\) satisfies the two symmetric relations in
\eqref{eq:fold-characteristic}.  Indeed, its sum and product are
\(\xi-\omega_L\) and \(-\omega_L\xi\), and the required identities
follow from
\[
 \mathfrak D_L=\frac{\zeta+\omega_L}{\zeta-\rho},\qquad
 \xi=\frac{\rho(\zeta+1)+\zeta(\omega_L-1)}{\zeta+\omega_L}.
\]
Positivity of \(\xi+\omega_L\) selects the distinct-variable branch of the
divided difference in \eqref{eq:fold-odd}.
The folded formula \eqref{eq:fold-odd} now gives the product directly.
The cross terms cancel because
\[
 x_*^2(\zeta-\xi)(\zeta+\omega_L)
 =\tau_*^2(1+\xi)(1-\omega_L),
\]
and \(U_{L-1}(-\omega_L)^2=(1-\omega_L^2)^{-1}\).  Simplification yields
\begin{equation}\label{eq:q-product-central}
\begin{aligned}
 &q_{2L-1}(x_*,\tau_*)q_{2L-1}(x_*,-\tau_*)\\
 &\quad=\frac{a^{2L-2}\tau_*^2}{(\xi+\omega_L)^2}
 \left(1-\frac{1-\omega_L}{\mathfrak D_L(\rho+\omega_L)}\right)
 \left[\frac{\mathfrak D_L(\rho+\omega_L)}{1+\omega_L}
       -(1+\xi)^2U_{L-1}(\xi)^2\right]>0.
\end{aligned}
\end{equation}
Indeed, \(\mathfrak D_L=S_b/Q_b>1\),
\(\rho+\omega_L>2\omega_L>1>1-\omega_L\), and the last bracket is
positive by \eqref{eq:central-target}.

At \(x=0\), the second-smallest singular value of \(A_{2L-1}\) is
\[
 s_{L,2}^{(0)}=\sqrt{1+a^2-2a\cos\frac{\pi}{L}}.
\]
Indeed, after separating odd and even coordinates, the two rectangular
off-diagonal blocks have Gram matrices, taken in their smaller dimensions, equal to the
\((L-1)\)-dimensional tridiagonal Toeplitz matrix with diagonal
\(1+a^2\) and off-diagonal \(a\).  Their singular values therefore have
squares
\[
 1+a^2+2a\cos\frac{k\pi}{L},\qquad 1\leq k\leq L-1,
\]
and each occurs twice, while the remaining singular value is zero.
We claim that
\begin{equation}\label{eq:x-c-second-singular}
 x_*+\tau_*<2\sqrt a\sin\frac{\pi}{2L}\leq s_{L,2}^{(0)}.
\end{equation}
Indeed, recall \(\rho_{\rm ref}=\cos\frac{\pi}{2L+1}\).  Since
\(\rho>\rho_{\rm ref}\) and \(U_{L-1}\) is strictly increasing above its
largest zero,
\[
 u_-=U_{L-1}(\rho)>U_{L-1}(\rho_{\rm ref})=:u_{-,\rm ref},
 \qquad u_{-,\rm ref}^{-2}=2(1-\rho_{\rm ref}).
\]
The identities \(\tau_*=a/u_-\) and \(\tau_*^2=2a(\zeta-\rho)\) imply
\(\zeta-\rho<a(1-\rho_{\rm ref})<1-\rho_{\rm ref}\).  Also
\[
 \frac{x_*^2}{2a}
 =\Delta_\omega\frac{\rho+\omega_L}{\zeta+\omega_L}
 <\Delta_\omega.
\]
Put
\[
 \theta_L=\frac{\pi}{4L},\qquad
 \Delta_L=\frac{\theta_L}{2L+1}.
\]
Then \(\Delta_\omega=2\sin^2\theta_L\) and
\(1-\rho_{\rm ref}=2\sin^2(\theta_L-\Delta_L)\).  Since all quantities are
positive, the preceding quotient bound for \(x_*^2\) gives
\[
 x_*<2\sqrt a\sin\theta_L.
\]
Likewise,
\[
 \tau_*^2=2a(\zeta-\rho)
 <2a^2(1-\rho_{\rm ref})
 =4a^2\sin^2(\theta_L-\Delta_L),
\]
and hence
\[
 \tau_*<2a\sin(\theta_L-\Delta_L)
 <2\sqrt a\sin(\theta_L-\Delta_L),
\]
where the last inequality uses \(0<a<1\), hence \(a<\sqrt a\).
It suffices to prove
\begin{equation}\label{eq:angle-inequality}
 \sin\theta_L+\sin(\theta_L-\Delta_L)<\sin(2\theta_L).
\end{equation}
The estimates
\[
 2\sin\theta_L-\sin(2\theta_L)
 =2\sin\theta_L(1-\cos\theta_L)<\theta_L^3
\]
and
\[
 \sin\theta_L-\sin(\theta_L-\Delta_L)
 =\int_{\theta_L-\Delta_L}^{\theta_L}\cos u\,du
 >\Delta_L\cos\theta_L
\]
reduce \eqref{eq:angle-inequality} to
\(\Delta_L\cos\theta_L>\theta_L^3\).  Substituting
\(\Delta_L=\theta_L/(2L+1)\) and \(\theta_L=\pi/(4L)\) gives the equivalent
condition \(2\cos\theta_L>\theta_L(\pi+2\theta_L)\).  For
\(0<\theta_L\leq\pi/8\), using \(\pi^2<10\),
\[
 2\cos\theta_L>2-\theta_L^2\geq\frac{59}{32}
 >\frac{25}{16}\geq\theta_L(\pi+2\theta_L).
\]
Consequently \(x_*+\tau_*<2\sqrt a\sin\frac{\pi}{2L}\).  Finally,
\[
 (s_{L,2}^{(0)})^2-4a\sin^2\frac{\pi}{2L}
 =(1-a)^2\geq0,
\]
which proves \eqref{eq:x-c-second-singular}.

Let \(\mu_1,\ldots,\mu_{2L-1}\) be the eigenvalues of the real symmetric
matrix
\[
 B_{2L-1}(x_*)=(x_*I_{2L-1}-A_{2L-1})J_{2L-1}.
\]
Right multiplication by \(J_{2L-1}\) preserves singular values, so the
numbers \(|\mu_i|\) are precisely the singular values of
\(x_*I_{2L-1}-A_{2L-1}\).  Moreover,
\[
 q_{2L-1}(x_*,\tau_*)q_{2L-1}(x_*,-\tau_*)
 =\prod_{i=1}^{2L-1}(\mu_i^2-\tau_*^2).
\]
Positivity of the product in \eqref{eq:q-product-central} implies that the
number of indices with \(|\mu_i|<\tau_*\) is even.  On the other hand, Weyl's
singular-value perturbation inequality and
\eqref{eq:x-c-second-singular} give, for the second-smallest singular value
\(s_{\rm next}\),
\[
 s_{\rm next}(x_*I_{2L-1}-A_{2L-1})
 \geq s_{L,2}^{(0)}-x_*>\tau_*,
\]
so at most one index can satisfy \(|\mu_i|<\tau_*\).  Parity therefore gives
\(|\mu_i|\geq\tau_*\) for every \(i\), and strict positivity of the product
rules out equality.  Hence
\(s_{\min}(x_*I_{2L-1}-A_{2L-1})>\mathfrak h_L^{\mathrm e}\).  Since
\(0<x_*<2r\sin\frac{\pi}{2L}\), the point lies strictly in the positive
central gap, and
\(\mathfrak h_{L-1}^{\mathrm o}>\mathfrak h_L^{\mathrm e}\).  Together
with the upper comparison
\(\mathfrak h_m^{\mathrm e}>\mathfrak h_m^{\mathrm o}\), this proves
\eqref{eq:central-interlace}.
\end{proof}

\subsection{Assembly of the global comparison}
\label{subsec:global-comparison}

\begin{proof}[Proof of \Cref{lem:mesh}]
Reflection and the central interlacing chain now assemble the local
comparisons into the global barrier inequality.  Put \(N=n+1\) and
\(M=N+1\), so the
larger and smaller matrices have orders \(M-1\) and \(N-1\), respectively.  If
\(\Sigma=\diag(1,-1,1,-1,\ldots)\), then
\(\Sigma A_n\Sigma=-A_n\), so reflection reduces every noncentral
comparison to a positive gap.  For each positive noncentral gap
\(G_{M-1,j}\), choose
\(x_j\in G_{M-1,j}^{\circ}\) attaining \(\gamma_{M-1,j}\).  Then
\(j\geq1\) and \((j+1)\pi/N\leq\pi/2\), so
\cref{lem:noncentral-gap-comparison} produces a point
\(\widehat x_j\in G_{N-1,j}^{\circ}\), and chord exhaustion gives
\[
 \gamma_{N-1,j}\geq g_{N-1}(\widehat x_j)
 >g_{M-1}(x_j)=\gamma_{M-1,j}.
\]
Reflection gives the same inequality on the corresponding reflected gap.

By \cref{lem:central-gap-comparison}, if \(M=2\ell_*+1\), the remaining
single central height of the larger, even-order matrix is
\(\mathfrak h_{\ell_*}^{\mathrm e}\),
whereas each central height of the
smaller odd-order matrix is
\(\mathfrak h_{\ell_*-1}^{\mathrm o}>\mathfrak h_{\ell_*}^{\mathrm e}\).  If
\(M=2\ell_*+2\), the two central heights of the larger odd-order matrix
equal
\(\mathfrak h_{\ell_*}^{\mathrm o}<\mathfrak h_{\ell_*}^{\mathrm e}\),
where \(\mathfrak h_{\ell_*}^{\mathrm e}\) is the single
central height of the smaller even-order matrix.
Thus every height of the larger problem is strictly smaller than at
least one height of the smaller problem.  By \eqref{eq:gamma-def}, choose
\(j_*\) with
\(\gamma_{M-1,j_*}=\max_j\gamma_{M-1,j}\).  Then
\[
 \gamma_{M-1}=\gamma_{M-1,j_*}
 <\max_k\gamma_{N-1,k}=\gamma_{N-1}.
\]
Since \(M-1=n+1\) and \(N-1=n\), this proves
\(\gamma_{n+1}<\gamma_n\).
\end{proof}



\section{Quantitative bounds and the transition scale}
\label{sec:topology-bounds}

The connectedness criterion \eqref{eq:connected-iff-gamma} and the strict
comparison of \cref{lem:mesh} reduce the remaining problem to estimating
\(\gamma_n\).  Put \(N=n+1\), and define
\begin{align}
 \mathcal U_n(a)&=r^n\csc\frac{\pi}{n+1},\label{eq:main-U-def}\\
 \mathcal L_n(a)&=
 \left(1+r^2-2r\cos\frac{\pi}{n+1}\right)
 \frac{r^n(1-r^2)}
 {(1-r^{2(n+1)})\sin\frac{3\pi}{2(n+1)}}.\label{eq:main-L-def}
\end{align}

\begin{proposition}[Optimal-order gap bounds]\label[proposition]{prop:gap-bounds}
For every \(n\geq2\),
\begin{equation}\label{eq:gamma-LU}
 \mathcal L_n(a)\leq\gamma_n\leq\mathcal U_n(a).
\end{equation}
In particular, both bounds have order \((n+1)a^{n/2}\), with constants
depending only on \(a\).
\end{proposition}

\begin{proof}
We prove the two inequalities in \eqref{eq:gamma-LU} separately.
For \(x=2r\cos\theta\in I_n^{\mathrm{sp}}\), set
\[
 \boldsymbol v=(v_j)_{j=1}^n,\qquad
 v_j=r^{j-1}\sin(j\theta),\qquad1\leq j\leq n.
\]
The residual \((xI_n-A_n)\boldsymbol v\) is supported in its last
coordinate, where it equals \(r^n\sin(N\theta)\).  Thus
\[
 g_n(2r\cos\theta)
 \leq\frac{r^n|\sin(N\theta)|}
 {\left(\sum_{j=1}^nr^{2j-2}\sin^2(j\theta)\right)^{1/2}}
 \leq r^n\frac{|\sin(N\theta)|}{\sin\theta}.
\]
Since \(\pi/N\leq\theta\leq(N-1)\pi/N\), this proves the upper bound.

For the lower bound, put
\(p_\theta=re^{\ii\theta}\), \(q_\theta=re^{-\ii\theta}\), so that
\(p_\theta+q_\theta=x\) and \(p_\theta q_\theta=a\).  For
\(\boldsymbol u=(u_j)_{j=1}^n\in\C^n\), put \(u_0=u_{n+1}=0\), and
define maps from \(\C^n\) to \(\C^{n+1}\) by
\[
 (\mathsf{C}_\theta\boldsymbol u)_j=p_\theta u_j-u_{j+1},\qquad
 (\mathsf{D}_\theta\boldsymbol u)_j=u_j-\overline{q_\theta}u_{j+1},
 \qquad0\leq j\leq n.
\]
Direct multiplication, using \(p_\theta+q_\theta=x\) and
\(p_\theta q_\theta=a\), gives
\begin{equation}\label{eq:rectangular-factorization}
 xI_n-A_n=\mathsf{D}_\theta^*\mathsf{C}_\theta.
\end{equation}
The two Gram matrices associated with the factorization
\eqref{eq:rectangular-factorization} are
diagonally unitarily similar to the tridiagonal Toeplitz matrix with
diagonal \(1+r^2\) and off-diagonal \(-r\).  Its eigenvalues are
\(1+r^2-2r\cos\frac{k\pi}{N}\), \(1\leq k\leq n\).  Consequently, for
every \(\boldsymbol u\in\C^n\),
\begin{equation}\label{eq:rectangular-min}
 \begin{aligned}
  \lVert\mathsf C_\theta\boldsymbol u\rVert_2
  &\geq \left(1+r^2-2r\cos\frac{\pi}{N}\right)^{1/2}
    \lVert\boldsymbol u\rVert_2,\\
  \lVert\mathsf D_\theta\boldsymbol u\rVert_2
  &\geq \left(1+r^2-2r\cos\frac{\pi}{N}\right)^{1/2}
    \lVert\boldsymbol u\rVert_2.
 \end{aligned}
\end{equation}
Put
\[
 \mathcal S_N=\sum_{j=0}^nr^{2j}=\frac{1-r^{2N}}{1-r^2}.
\]
Unit normal vectors to the two hyperplane ranges are
\[
 \boldsymbol u_C=\mathcal S_N^{-1/2}
 (\overline{p_\theta}^{\,n},\overline{p_\theta}^{\,n-1},\ldots,1)^T,
 \qquad
 \boldsymbol u_D=\mathcal S_N^{-1/2}(1,q_\theta,\ldots,q_\theta^n)^T.
\]
The geometric sum gives
\begin{align}
 |\langle\boldsymbol u_D,\boldsymbol u_C\rangle|
 &=\frac{r^n}{\mathcal S_N}
 \left|\sum_{j=0}^{N-1}e^{\ii(2j-(N-1))\theta}\right|\notag\\
 &=
 \frac{r^n(1-r^2)}{1-r^{2N}}
 \left|\frac{\sin(N\theta)}{\sin\theta}\right|
 =:\chi_N(\theta).
\label{eq:chi}
\end{align}
Let \(P_D\) be the orthogonal projection onto the range of
\(\mathsf D_\theta\).  If \(\boldsymbol w\perp\boldsymbol u_C\),
Cauchy--Schwarz applied to the component of \(\boldsymbol u_D\)
orthogonal to \(\boldsymbol u_C\) gives
\[
 |\langle\boldsymbol u_D,\boldsymbol w\rangle|^2
 \leq(1-\chi_N(\theta)^2)\lVert\boldsymbol w\rVert_2^2,
 \qquad
 \lVert P_D\boldsymbol w\rVert_2\geq
 \chi_N(\theta)\lVert\boldsymbol w\rVert_2.
\]
Put \(\kappa_N=1+r^2-2r\cos(\pi/N)\).
By \eqref{eq:rectangular-min}, the restriction of
\(\mathsf D_\theta^*\) to the range of \(\mathsf D_\theta\) has least
singular value at least \(\sqrt{\kappa_N}\).  Hence
\[
 \lVert\mathsf D_\theta^*\mathsf C_\theta\boldsymbol u\rVert_2
 \geq\sqrt{\kappa_N}\lVert P_D\mathsf C_\theta\boldsymbol u\rVert_2
 \geq\kappa_N\chi_N(\theta)\lVert\boldsymbol u\rVert_2.
\]
Together with \eqref{eq:rectangular-factorization}, this gives
\[
 g_n(2r\cos\theta)\geq
 \left(1+r^2-2r\cos\frac{\pi}{N}\right)\chi_N(\theta).
\]
Putting \(\theta=3\pi/(2N)\), which lies in the first spectral gap,
proves the lower bound.

Finally,
\[
 \mathcal U_n(a)\leq\frac12(n+1)r^n,
\]
while
\begin{equation}\label{eq:L-simple}
 \mathcal L_n(a)\geq C_{\rm lb}(a)(n+1)r^n,\qquad
 C_{\rm lb}(a)=\frac{2(1-r)^2(1-r^2)}{3\pi}.
\end{equation}
The estimate \eqref{eq:L-simple} and the preceding upper estimate follow from
\(\sin\frac{\pi}{N}\geq2/N\),
\(1+r^2-2r\cos\frac{\pi}{N}\geq(1-r)^2\), and
\(\sin\frac{3\pi}{2N}\leq3\pi/(2N)\).  This proves the asserted optimal
order up to constants depending only on \(a\).
\end{proof}

The bounds also locate the first connected order:
\begin{equation}\label{eq:N-two-sided-main}
 \begin{aligned}
  1+\max\bigl(\{1\}\cup
  \{n\geq2:\mathcal L_n(a)\geq\varepsilon\}\bigr)
  &\leq N_c(a,\varepsilon),\\
  N_c(a,\varepsilon)
  &\leq\min\{n\geq2:\mathcal U_n(a)<\varepsilon\}.
 \end{aligned}
\end{equation}

\begin{proof}[Proof of \Cref{thm:canonical-main}]
By \Cref{prop:gap-bounds},
\(\gamma_n\leq\mathcal U_n(a)\to0\).  \Cref{lem:mesh} says that
\(\gamma_n\) is strictly decreasing.  Equation
\eqref{eq:connected-iff-gamma} therefore shows that, once the
pseudospectrum is connected, it stays connected at every larger matrix
order.  Contractibility and the connectedness criterion were proved in
Proposition~\ref{prop:topology}.

If \(\mathcal L_n(a)\geq\varepsilon\), then
\(\gamma_n\geq\varepsilon\), so the pseudospectrum at order \(n\)
is disconnected.  If \(\mathcal U_n(a)<\varepsilon\), then
\(\gamma_n<\varepsilon\), so it is connected.  Strict decrease of
\(\gamma_n\) now gives \eqref{eq:N-two-sided-main}.

Put \(N=N_c(a,\varepsilon)\), \(\Lambda=\log(1/\varepsilon)\), and
\(\beta=-\log r=\tfrac12\log(1/a)>0\).
For fixed \(a\), positivity of every \(\gamma_n\) implies
\(N\to\infty\) as \(\varepsilon\downarrow0\).  Thus \(N\geq3\) for
all sufficiently small \(\varepsilon\), and minimality of \(N\) gives
\[
 C_{\rm lb}(a)(N+1)e^{-\beta N}
 \leq\gamma_N<\varepsilon\leq\gamma_{N-1}
 \leq\tfrac12Ne^{-\beta(N-1)}.
\]
Taking logarithms yields
\begin{equation}\label{eq:critical-log-sandwich}
 \Lambda+\log(N+1)+\log C_{\rm lb}(a)
 <\beta N\leq\Lambda+\log N+\beta-\log2.
\end{equation}
The lower inequality bounds \(N\) below by a positive multiple of
\(\Lambda\) for large \(\Lambda\).  In the upper inequality,
\(\log N\leq\beta N/2\) for large \(N\), which gives the corresponding
upper bound.  Hence \(\log N=\log\Lambda+O_a(1)\), and substitution into
\eqref{eq:critical-log-sandwich} proves \eqref{eq:N-asymptotic-main}.

The sole gap at order two is central, so the central-height formula gives
\(\gamma_2=\mathfrak h_1^{\mathrm e}=s_{\min}([a])=a\).
Equation~\eqref{eq:connected-iff-gamma} therefore gives
\(N_c(a,\varepsilon)=2\) if and only if \(\varepsilon>a\).
This proves all assertions.
\end{proof}



\section{The boundary cases}
\label{sec:boundaries}

The same matrices at \(a=0\) and \(a=1\) have simpler pseudospectra.

\begin{proposition}\label[proposition]{prop:boundary-cases}
Let \(n\geq2\) and \(\varepsilon>0\).
The pseudospectrum of \(A_n(0)\) is an open disk centered at zero.
For \(a=1\), every pseudospectral component is contractible, and the
pseudospectrum is connected exactly when \(\varepsilon>\eta_n\), where
\begin{equation}\label{eq:normal-threshold}
 \eta_n=\begin{cases}
 2\sin\dfrac{\pi}{2(n+1)},&n\text{ even},\\[5pt]
 \sin\dfrac{\pi}{n+1},&n\text{ odd}.
 \end{cases}
\end{equation}
Moreover, \(\eta_{n+1}<\eta_n\) for all \(n\geq2\), and
\(\eta_n=\pi/(n+1)+O(n^{-3})\).
\end{proposition}

\begin{proof}
At \(a=0\), the matrix is the nilpotent shift \(V_n\).
The proof of Lemma~\ref{lem:component} applies unchanged, so every
pseudospectral component contains its only eigenvalue, zero; the set is
therefore connected.  For \(|q|=1\), the unitary matrix
\(D(q)=\diag(1,q,\ldots,q^{n-1})\) satisfies \(D(q)^*V_nD(q)=qV_n\).
The pseudospectrum is consequently invariant under all rotations.  Its
set of moduli is a connected interval containing zero; boundedness and
openness make it \([0,R)\) for some \(R>0\).  Hence it is an open disk.

At \(a=1\), the matrix is real symmetric, and its pseudospectrum is the
union of open disks of radius \(\varepsilon\) centered at
\(2\cos(k\pi/(n+1))\), \(1\leq k\leq n\).  It is connected precisely
when twice the radius exceeds every adjacent eigenvalue gap.  The identity
\[
 \cos\frac{k\pi}{n+1}-\cos\frac{(k+1)\pi}{n+1}
 =2\sin\frac{(2k+1)\pi}{2(n+1)}\sin\frac{\pi}{2(n+1)}
\]
gives \eqref{eq:normal-threshold}.  Since contraction of the imaginary
part decreases the distance to every eigenvalue, the retraction argument
of Proposition~\ref{prop:topology} proves contractibility of each component.

For the transition from \(n=2m\) to \(2m+1\), set
\(x=\pi/(4m+2)\) and \(y=\pi/(2m+2)\).  Then
\(y<2x<\pi/2\), so \(\sin y<\sin(2x)<2\sin x\).
For the transition from \(2m+1\) to \(2m+2\), put
\(p=m+1\geq2\), \(x=\pi/(2p)\), and \(y=\pi/(4p+2)\).
The inequalities \(\sin x>x-x^3/6\), \(\sin y<y\), and
\(24p^2>\pi^2(2p+1)\) give \(\sin x>2\sin y\).
The last inequality holds because \(24p^2/(2p+1)\) increases for
\(p\geq2\) and equals \(96/5>\pi^2\) at \(p=2\).
Finally, applying \(0\leq t-\sin t\leq t^3/6\) to each parity formula yields
\[
 0\leq\frac{\pi}{n+1}-\eta_n\leq\frac{\pi^3}{6(n+1)^3},
\]
which proves the stated asymptotic.
\end{proof}

Thus the connectedness threshold is zero at \(a=0\), has order
\((n+1)a^{n/2}\) for fixed \(0<a<1\), and has order \((n+1)^{-1}\)
at \(a=1\).


\section{Numerical illustrations}\label{sec:numerics}

\Cref{fig:transition,fig:barriers} illustrate two aspects of the matrix-order
transition: coalescence of pseudospectral components and decay of the
real-axis backward-error barrier.

At each sampled point, the least singular value was computed by a
double-precision singular value decomposition.  The contour plots use a
\(561\)-by-\(361\) Cartesian grid.  Each spectral gap was scanned at \(121\)
equally spaced points before bounded scalar optimization was used to estimate
\(\gamma_n\).  The displayed component counts were also evaluated on
\(1001\) points per gap, with the 1-Lipschitz dependence of the least singular
value controlling variation between adjacent sample points.

\begin{figure}[!htbp]
 \centering
 \includegraphics[width=\textwidth]{fig1_connectedness_transition}
 \caption{Numerically sampled open pseudospectra of \(A_n(1/4)\) at
 \(\varepsilon=10^{-2}\).  Shading marks the sampled pseudospectrum and
 crosses mark the exact eigenvalues.  The samples show component
 coalescence between \(n=6\) and \(n=7\).}
 \label{fig:transition}
\end{figure}

\begin{figure}[!htbp]
 \centering
 \includegraphics[width=.84\textwidth]{fig2_gap_barrier_bounds}
 \caption{Real-axis backward-error barriers for \(a=1/4\).
 The gray band, bounded by the dotted lower curve \(\mathcal L_n(a)\) and
 dashed upper curve \(\mathcal U_n(a)\), is the rigorous enclosure of
 \(\gamma_n\) and establishes the order \((n+1)a^{n/2}\).  Blue circles joined by
 a solid line are floating-point estimates; the red horizontal line is the
 fixed pseudospectral level \(\varepsilon=10^{-2}\), and the black vertical line marks
 the numerical crossing \(\widehat N_c=7\).  The theorem gives the bracket
 \(6\leq N_c\leq9\).}
 \label{fig:barriers}
\end{figure}




\FloatBarrier

\section{Concluding remarks}

For \(A_n(a)\), contraction toward the real axis reduces connectedness to
the maximum of one real function over the spectral interval.  The signed
pencil then permits a strict comparison of these maxima at consecutive
orders, despite the change in central-gap geometry.  The bounds in
Proposition~\ref{prop:gap-bounds} identify the polynomial factor multiplying
the exponential decay and yield the second logarithmic term in the
first connected order.

The asymptotic formula is for fixed \(0<a<1\).  A joint limit in which
\(a\) approaches one as \(n\) increases would have to connect the
exponential scale of Section~\ref{sec:topology-bounds} with the algebraic
scale of Section~\ref{sec:boundaries}.

\funding{The author received no specific funding for this work.}
\dataavailability{The manuscript sources and code for the numerical illustrations
 are available at \url{https://github.com/jinshanmu/Hatano-Nelson},
 in the \texttt{preprint} directory.}
\acknowledgments{OpenAI Codex assisted with manuscript restructuring,
 language editing, algebraic checks, and verification of bibliographic metadata.}
\conflictsofinterest{The author declares no conflict of interest.}

\reftitle{References}
\bibliography{references}
\end{document}
